% DES-LOC: Desynced Low Communication Adaptive Optimizers for Foundation Models
% Reconstructed from 12473_DES_LOC_Desynced_Low_Com.pdf
% Claude-1 (M001-M025): Sections 1-7 main body LaTeX reconstruction
% Claude-2 (M026-M050): Section 3 Convergence Guarantees — strict 1:1 PDF match, ψ formula fix
% Claude-3 (M051-M075): Section 4 Experimental Design — Metrics paragraph completion, Section 5.2 fix
% Claude-4 (M076-M100): Section 5.1-5.4 1:1 PDF校验 — Fig5 caption, Table 1/2 caption+data补全, Section 5.4段落补全


\documentclass{article}
\usepackage{neurips_2026}
\usepackage{amsmath,amssymb,amsfonts}
\usepackage{algorithmic}
\usepackage{algorithm}
\usepackage{graphicx}
\usepackage{hyperref}
\usepackage{booktabs}
\usepackage{multirow}

\title{DES-LOC: Desynced Low Communication Adaptive Optimizers for Foundation Models}

\author{
Alex Iacob$^{\dagger,1,2}$ \quad Lorenzo Sani$^{1,2}$ \quad Mher Safaryan$^{3}$ \quad Andrej Jovanovi\'{c}$^{*,1}$ \\
William F.\ Shen$^{1}$ \quad Xinchi Qiu$^{1}$ \quad Samuel Horv\'{a}th$^{*,5}$ \\
Preslav Aleksandrov$^{1,6}$ \quad Paris Giampouras$^{*,4}$ \quad Meghdad Kurmanji$^{*,1}$ \\
Nicholas D.\ Lane$^{1,2}$ \\
\\
$^{1}$ University of Cambridge \quad $^{2}$ Flower Labs \quad $^{3}$ ISTA \\
$^{4}$ University of Warwick \quad $^{5}$ MBZUAI \quad $^{6}$ IMEC \\
$^{\dagger}$\texttt{aai30@cam.ac.uk}; $^{*}$ Equal contributions
}

\begin{document}
\maketitle

\begin{abstract}
Scaling foundation model training with Distributed Data Parallel (DDP) methods
is bandwidth-limited. Existing infrequent communication methods like Local
SGD were designed to synchronize model parameters only and cannot be trivially
applied to adaptive optimizers due to additional optimizer states. Heuristic approaches that keep states local or reset them lack guarantees and can be unstable in
compute-efficient batch regimes; conversely, Local Adam synchronizes all states
uniformly and is provably convergent but triples communication costs. We propose
Desynced Low Communication Adaptive Optimizers (DES-LOC), a family of optimizers assigning independent synchronization periods to parameters and momenta,
enabling lower communication costs while preserving convergence. Our theoretical
analysis shows that while parameter synchronization dominates the asymptotic rate
in-expectation, high-probability convergence guarantees require at least infrequent
synchronization of the second momentum. Furthermore, we prove that more frequent momentum sync permits larger stable step sizes. Experiments on language
models of up to 1.7B show that DES-LOC can communicate $170\times$ less than DDP
and $2\times$ less than the previous state-of-the-art Local Adam, enabling $1.3$--$2.1\times$
wall-clock speedups over DDP for 1--13B models on 100Gb/s links. Furthermore,
unlike previous heuristic methods, DES-LOC is robust to worker failures offering
a scalable, efficient, and fault-tolerant solution for foundation model training.
\end{abstract}

%=============================================================================
\section{Introduction}
\label{sec:introduction}
%=============================================================================

Training foundation models requires distributing optimization across workers for improved memory
and compute. However, frequent gradient communication in standard Distributed Data Parallelism
(DDP)~\citep{li2020pytorch} increases networking costs and limits scalability. Early works like Local
SGD~\citep{stich2019local} and FedAvg~\citep{mcmahan2017communication} reduced this overhead by synchronizing
infrequently, averaging parameters only after $K \gg 1$ local steps, instead of gradients at every step.
However, modern foundation model training, e.g., Large Language Models~\citep{dubey2024llama}, uses
adaptive optimizers~\citep{kingma2015adam} which require additional momenta.

Some extensions of Local SGD to adaptive optimizers~\citep{sani2025photon,douillard2023diloco}
only average model parameters, which poses challenges. First, they lack convergence guarantees.
Second, keeping momenta local~\citep{douillard2023diloco} accumulates noisy small-batch gradients and
provides no means to initialize workers. This makes them unsuitable for failure-prone environments.
Third, re-initializing momenta~\citep{sani2024future,sani2025photon} destabilizes training.

Local Adam~\citep{cheng2025convergence} addresses these challenges, proving periodic synchronization
can converge faster than standard Adam with DDP, and remain robust to the addition of
new workers. However, it requires synchronizing momenta alongside model parameters, tripling
communication costs compared to Local SGD. Hence, we aim to answer the following question:

\begin{center}
\emph{Can independently syncing parameters and momenta improve communication
efficiency for adaptive optimizers while maintaining convergence and robustness?}
\end{center}

As a result of our inquiry, we propose a new optimizer family, Desynced Low Communication
Adaptive Optimizers (DES-LOC), which sets independent synchronization frequencies for parameters
and momenta. This reduces communication overhead by synchronizing momenta less frequently.

\paragraph{Contributions:}
\begin{enumerate}
    \item \textbf{Provable convergence.} We prove convergence (see Section~\ref{sec:convergence}) for DES-LOC under: non-convex
    objectives when using SGD with momentum (SGDM), and weakly convex objectives when using Adam. Our theory indicates a higher momentum sync frequency enables larger step sizes.
    Furthermore, high-probability bounds demand momenta be synced with finite period for $\beta_2 < 1.0$.

    \item \textbf{Communication reduction.} We empirically show that parameters require more frequent sync than
    momenta, and that less frequent momentum sync reduces communication costs ($2\times$ vs Local
    Adam, $170\times$ vs DDP), leading to $1.3$--$2.1\times$ reductions in training time over DDP on our hardware.

    \item \textbf{Scalability to large models.} We validate DES-LOC at billion-scale language model training,
    demonstrating competitive ICL performance against both Local Adam and DDP.

    \item \textbf{Hardware robustness.} Unlike previous heuristic methods, DES-LOC avoids persistent local states,
    enabling it to seamlessly integrate new workers to support environments prone to system failures.
\end{enumerate}

%=============================================================================
\section{Desynced Low Communication Adaptive Optimizers (DES-LOC)}
\label{sec:desloc}
%=============================================================================

We start by characterizing the relation between the rate of change of optimizer states and Local
Adam, and how these can be leveraged to lower the communication cost. Consider the Adam update:
\begin{equation}
u_t = \beta_1 u_{t-1} + (1 - \beta_1) g_t \quad \text{and} \quad v_t = \beta_2 v_{t-1} + (1 - \beta_2) g_t \odot g_t.
\end{equation}

For Local Adam, convergence is contingent on $\beta_2$ satisfying $1 - \beta_2 = \tilde{O}(K^{-3/2} R^{-1/2})$~\citep{cheng2025convergence} where $K$ is the number of local steps and $R$ the total communication rounds. Large
$K$ or $R$ implies $\beta_2 \to 1$, and conversely larger $\beta_2$ permits higher $K$ or $R$.

A useful summary measure is the number of steps until a state's weight decays to a fraction $\psi$,
$\tau_\psi(\beta) = \frac{\ln \psi}{\ln \beta}$. Following~\citet{pagliardini2025ademamix}, we use the half-life $\tau_{0.5}$ as our primary measure,
omitting $\beta$ when clear. For typical values of $\beta$, we have $\tau_{0.5}(0.95) \approx 13.5$~\citep{allal2025smollm2},
$\tau_{0.5}(0.999) \approx 692.8$~\citep{kingma2015adam}, and $\tau_{0.5}(0.9999) \approx 6931$~\citep{taniguchi2024adopt}.

Intuitively, larger half-lives imply synchronizing gradients over longer horizons as the optimizer
is less sensitive to new gradients; choosing $\beta = 0$ ignores all previous momenta, whereas $\beta \to 1$
progressively attenuates signal from the current gradient.

While the half-life captures the horizon for which an optimizer state remains relevant to model
updates, it provides no information on its absolute rate of change. With coordinate-wise clipping,
each gradient component satisfies $|(g_t)_i| \leq \rho$. Unrolling Adam's recursions over $K$ local steps gives
the follow relation: $u_{t+K} = \beta_1^K u_t + (1 - \beta_1) \sum_{k=0}^{K-1} \beta_1^k g_{t+K-1-k}$ and its second moment analogue.
Since $|g_{t,i}| \leq \rho$ and $|(g_t \odot g_t)_i| \leq \rho^2$, the maximal $\ell_\infty$ drift of each moment is (see Section~\ref{sec:drift_derivation}):
\begin{align}
\|u_{t+K} - u_t\|_\infty &\leq 2\rho(1 - \beta_1^K), \label{eq:drift_u} \\
\|v_{t+K} - v_t\|_\infty &\leq 2\rho^2(1 - \beta_2^K). \label{eq:drift_v}
\end{align}

From the above, large $\beta$ values and small clip bounds $\rho$, a common practice in foundation model
training~\citep{brown2020language,scao2022bloom}, limit the absolute changes in optimizer states. We can
construct similar reasoning for other optimizers~\citep{sutskever2013importance,taniguchi2024adopt}, and
norm-based clipping~\citep{pascanu2013difficulty,brown2020language}. From the above, the half-life of an
optimizer state should inform its synchronization frequency. For example, if $\tau_{0.5}(0.95) \approx 13.5$ and
$K = 256$, synchronization only affects few initial local steps. Over the course of the local training,
the impact of the synchronised optimizer state shall decay to $0$ given Equations~\ref{eq:drift_u} and~\ref{eq:drift_v}. Conversely,
if $K = 16$, synchronization approximately matches the half-life, strongly influencing local updates.

%-----------------------------------------------------------------------------
\subsection{DES-LOC Algorithm}
\label{sec:desloc_algorithm}
%-----------------------------------------------------------------------------

\begin{algorithm}[t]
\caption{DES-LOC}
\label{alg:desloc}
\begin{algorithmic}[1]
\REQUIRE Model tensors, update functions, hyper-parameters
\STATE $x_0 \in \mathbb{R}^d$, $\{s^j_{-1}\}_{j=1}^N \in (\mathbb{R}^d)^N$ \COMMENT{initial parameter vector, the initial $N$ optimizer states}
\STATE $\{\text{UPDATE}^j\}_{j=1}^N : (\mathbb{R}^d \times \mathbb{R}^d \to \mathbb{R}^d)^N$ \COMMENT{updates optimizer state $j$ from its previous state and the gradient}
\STATE $\text{OPT} : \mathbb{R}^d \times \mathbb{R}^d \times \mathbb{R}_+ \times (\mathbb{R}^d)^N \to \mathbb{R}^d$ \COMMENT{update params from all worker models}
\STATE $\text{SERVEROPT} : \mathbb{R}^d \to \mathbb{R}^d$ \COMMENT{update params using an abstract outer optimizer}
\STATE $\rho \in \mathbb{R}_+$, $\{\eta_t\}_{t=0}^{T-1} \in (\mathbb{R}_+)^{T-1}$ \COMMENT{clipping radius, learning-rate schedule}
\STATE $T, M \in \mathbb{N}_+$ \COMMENT{total optimization steps and number of workers}
\STATE $K_x \in \mathbb{N}_+$, $\{K_j\}_{j=1}^N \in (\mathbb{N}_+)^N$ \COMMENT{communication periods (steps)}
\ENSURE $x_T$, $\{s^j_{T-1}\}_{j=1}^N$
\STATE \textbf{for each worker} $m$: $x^m_0 \gets x_0$, $s^{j,m}_{-1} \gets s^j_{-1}$ \COMMENT{local init}
\FOR{$t = 0, \ldots, T-1$} \COMMENT{training loop}
    \FOR{all workers $m = 0, \ldots, M-1$ \textbf{in parallel}}
        \STATE $g^m_t \gets \nabla F(x^m_t; \xi^m_t)$ \COMMENT{stochastic grad}
        \STATE $\hat{g}^m_t \gets \text{clip}(g^m_t, \rho)$ \COMMENT{per-coordinate clipping}
        \FOR{$j = 1$ \TO $N$}
            \IF{$t \bmod K_j = 0$} \COMMENT{sync $s^j$}
                \STATE $s^{j,m}_t \gets \text{UPDATE}^j\left(\mathbb{E}_m[s^{j}_{t-1}], \hat{g}^m_t\right)$
            \ELSE
                \STATE $s^{j,m}_t \gets \text{UPDATE}^j\left(s^{j,m}_{t-1}, \hat{g}^m_t\right)$
            \ENDIF
        \ENDFOR
        \IF{$t \bmod K_x = 0$} \COMMENT{sync $x$}
            \STATE $x^m_{t+1} \gets \text{OPT}\left(\text{SERVEROPT}(\mathbb{E}_m[x^m_t]), \hat{g}^m_t, \eta_t, \{s^{j,m}_t\}_{j=1}^N\right)$
        \ELSE
            \STATE $x^m_{t+1} \gets \text{OPT}\left(x^m_t, \hat{g}^m_t, \eta_t, \{s^{j,m}_t\}_{j=1}^N\right)$
        \ENDIF
    \ENDFOR
\ENDFOR
\end{algorithmic}
\end{algorithm}

Motivated by the above insights, we formalize Desynced Low Communication Adaptive Optimizers
as a family of optimizers offering the same convergence and robustness as Local Adam but with
significantly lower communication costs. Our approach applies generically to adaptive optimizers
parameterized by $\text{OPT} : (\mathbb{R}^d, \mathbb{R}^d, \mathbb{R}_{>0}, \{\mathbb{R}^d\}^N) \to \mathbb{R}^d$, with $N$ optimizer states $\{s^j_{-1}\}_{j=1}^N \subset \mathbb{R}^d$,
each updated by $\text{UPDATE}^j : (\mathbb{R}^d, \mathbb{R}^d) \to \mathbb{R}^d$. Coordinate-wise clipping is defined as $[\text{clip}(X, \rho)]_i = \text{sgn}(X_i) \cdot \min\{|X_i|, \rho\}$. To ensure that our method is provably convergent, SERVEROPT is that of
FedAvg~\citep{mcmahan2017communication}. However, our algorithm directly extends to the larger FedOpt~\citep{reddi2021adaptive} framework, which we discuss in Section~\ref{sec:fedopt}.

We focus our analysis on SGDM and Adam. As shown in Algorithm~\ref{alg:desloc}, DES-LOC synchronizes
parameters $x \in \mathbb{R}^d$ and optimizer states $\{s^j\}_{j=1}^N$ at state-specific intervals $K_x, \{K_j\}_{j=1}^N \in \mathbb{N}_+$.
Setting $N = 2$, $s^1_t = u_t$, $s^2_t = v_t$, and using update rules $\text{UPDATE}^1$, $\text{UPDATE}^2$ based on the Adam
update rules above yields DES-LOC-Adam (see Algorithm~\ref{alg:desloc_adam}).

\paragraph{Toy Example} To highlight DES-LOC's practical benefit, Fig.~\ref{fig:toy} illustrates a scenario where DES-LOC
and Local Adam converge under noisy gradients, while prior heuristic methods~\citep{douillard2023diloco,sani2025photon,iacob2025dept,sani2024future} fail.

%=============================================================================
\section{Convergence Guarantees for DES-LOC}
\label{sec:convergence}
%=============================================================================
% Claude-2 (M026-M050): Strict 1:1 with PDF pages 3-5. ψ formula matches PDF Eq.(4).

This section provides theoretical support for the proposed DES-LOC approach. We focus on a version
of the Adam optimizer that uses only a single momentum state. Extensions to the full Adam optimizer
with both momenta are available in Section~\ref{sec:adam_extension} with high-probability bounds shown in Section~\ref{sec:high_prob}.

Formally, we consider the following optimization problem:
\begin{equation}
\min_{x \in \mathbb{R}^d} f(x) := \frac{1}{M} \sum_{m=1}^{M} f_m(x), \quad \text{with } f_m(x) = \mathbb{E}_{\xi \sim \mathcal{D}_m}[F_m(x; \xi)].
\label{eq:objective}
\end{equation}

In this setup, all $M$ machines collaboratively minimize the objective in~\eqref{eq:objective}. Generally, we assume
each machine $m$ has access to only dataset $\mathcal{D}_m$, which can differ from device to device. This recovers
the homogeneous distribution case when all machines have the same dataset $\mathcal{D}_1 = \mathcal{D}_2 = \cdots = \mathcal{D}_M$
and minimize the same loss $f_1(x) = f_2(x) = \cdots = f_m(x) = f(x)$. We assume each machine $m$
computes mini-batch stochastic gradients corresponding to randomly selected samples $\xi \sim \mathcal{D}_m$ from
dataset $\mathcal{D}_m$. We further use the following technical assumptions on the problem structure.

% Figure 1: PDF page 4 top — placed here to match PDF layout (between Section 2 Toy Example and Assumptions)
\begin{figure}[t]
\centering
% [Figure 1 graphic from original PDF]
\caption{We present: (left) the distance to the optimum and (right) a 2-D contour of a toy problem
where DES-LOC ($K_x = 192$, $K_u = 192$, $K_v = 692$) and Local Adam ($K = K_x$) both converge to the
optimum (overlapping). Methods keeping optimizer states local (\protect\rule{6pt}{6pt}) fail to converge.
Periodically resetting states (\protect\rule{6pt}{6pt}) similarly stalls due to repeated oscillations.
We optimize the non-convex function $f(x_1, x_2) = (1-x_1)^2 + 100(x_2 - x_1^2)^2$ with $M = 256$
workers and IID Gaussian noise ($\sigma = 1.5$). We show an example of such a toy problem on Non-IID
data in Fig.~8.}
\label{fig:toy}
\end{figure}

\begin{assumption}[Lower bound and smoothness]
\label{asm:smooth}
The overall loss function $f: \mathbb{R}^d \to \mathbb{R}$ is lower bounded by some $f^* \in \mathbb{R}$ and all local loss functions $f_m$ are $L$-smooth:
\begin{equation}
\|\nabla f_m(x) - \nabla f_m(y)\| \leq L\|x - y\|, \quad \text{for any } x, y \in \mathbb{R}^d.
\end{equation}
\end{assumption}

\begin{assumption}[Unbiased noise with bounded stochastic variance]
\label{asm:noise}
The stochastic gradient $g^m$ of local loss function $f_m$ computed by machine $m$ is unbiased and the noise has bounded variance:
\begin{equation}
\mathbb{E}[g^m] = \nabla f_m(x), \quad \mathbb{E}[\|g^m_t - \nabla f_m(x)\|^2] \leq \sigma^2, \quad \text{for any } x \in \mathbb{R}^d.
\end{equation}
\end{assumption}

\begin{assumption}[Bounded heterogeneity]
\label{asm:heterogeneity}
For any $x \in \mathbb{R}^d$, the heterogeneity is bounded by
\begin{equation}
\frac{1}{M} \sum_{m=1}^{M} \|\nabla f_m(x)\|^2 \leq G^2 + B^2 \|\nabla f(x)\|^2.
\end{equation}
\end{assumption}

All three assumptions are standard and widely used in the convergence analysis of optimization
algorithms~\citep{yu2019linear,karimireddy2020scaffold,wang2021field,yuan2022revisiting}.
Note that the bounded heterogeneity condition recovers the homogeneous case when $G^2 = 0$ and
$B^2 = 1$. To facilitate the technical presentation of the analysis, we view model and optimizer
state synchronizations through assigning probabilities to each averaging event. Particularly, instead
of averaging model parameters every $K_x$ steps (i.e., $t \bmod K_x = 0$), we average with probability
$p_x = \frac{1}{K_x}$, which are statistically equivalent. In the following theorem, we provide convergence rate
of SGDM optimizer under such probabilistic and decoupled synchronization:

\begin{theorem}
\label{thm:sgdm}
Let Assumptions~\ref{asm:smooth},~\ref{asm:noise} and~\ref{asm:heterogeneity} hold. Then, choosing the step size $\eta = \min(\eta_0, \frac{1}{\sqrt{T}})$ with
\begin{equation}
%%% BUG-1 FIX: 因子6在η₀中(1/(6√…)), 不在ψ中. 严格匹配PDF page 4 Eq.(4).
\eta_0 \stackrel{\text{def}}{=} \frac{1}{4L} \min\left(1 - \beta, \frac{1}{6\sqrt{\psi \max(1, B^2 - 1)}}\right), \quad \text{where } \psi \stackrel{\text{def}}{=} \frac{4(1-p_x)}{p_x^2} \cdot \frac{(1-\beta)(1-p_u)}{1-(1-p_u)\beta},
\label{eq:step_size}
\end{equation}
the average iterates $\bar{x}_t = \mathbb{E}_m[x^m_t]$ of DES-LOC-SGDM converge with the following rate:
\begin{equation}
\frac{1}{T} \sum_{t=0}^{T-1} \mathbb{E}\|\nabla f(\bar{x}_t)\|^2 \leq \frac{4}{\sqrt{T}} \left(f(x_0) - f^* + \frac{L\sigma^2}{2M}\right) + \mathcal{O}\left(\frac{1+\psi}{T}\right).
\label{eq:rate}
\end{equation}
\end{theorem}

We now discuss the convergence result and its implications. The obtained rate~\eqref{eq:rate} is asymptotically
optimal for this setup~\citep{arjevani2023lower}. Notably, the leading term $\mathcal{O}(\frac{1}{\sqrt{T}})$ is unaffected by the
number of local steps. Interestingly, probabilities $p_x$, $p_u$, and the momentum parameter $\beta$ appear in
the higher-order term $\mathcal{O}(\frac{1}{T})$, and thus have a limited impact on asymptotic convergence speed.

Regarding state synchronization, it is evident from~\eqref{eq:step_size} that model synchronization has a greater
impact on convergence due to the dependence $\psi = \mathcal{O}(\frac{1}{p_x^2})$. With vanishing $p_x$, the $\psi$ term becomes
unbounded and breaks the rate. For optimizer states, it seems that momentum averaging can be
turned off ($p_u = 0$) without affecting the asymptotic behavior of the rate. Setting $p_x = 1$ and $p_u = 0$
recovers standard mini-batch SGDM~\citep{liu2020improved}. However, the $\psi$ term also appears in the step-size restriction~\eqref{eq:step_size}. As $p_u \to 0$, $\frac{(1-p_u)}{1-(1-p_u)\beta} \to \frac{1}{1-\beta}$. This imposes the most severe restriction on the
learning rate $\eta_0$ since $\eta_0 \propto \frac{1}{\sqrt{\psi}}$, as $\psi$ is maximized. This theory shows that increasing the frequency
$p_u$ of momentum averaging---while not changing the asymptotic rate---allows for a larger step size,
potentially leading to faster convergence in practice. This theory justifies that momentum states
can be synchronized less frequently than parameters and that more averaging improves convergence
by supporting larger step sizes. Furthermore, our high probability analysis of DES-LOC-Adam in
Section~\ref{sec:high_prob} shows that the sync frequency of momenta must be finite for $\beta_2 < 1.0$.

%=============================================================================
\section{Experimental Design}
\label{sec:experimental_design}
%=============================================================================

Our experimental setup addresses the following research questions:
\begin{description}
\item[RQ1] Do theoretical rates of change predict the empirical evolution of optimizer states?
\item[RQ2] How does the synchronization frequency of a model/optimizer state impact performance?
\item[RQ3] To what extent can DES-LOC cut communication w.r.t.\ Local Adam in practical scenarios?
\item[RQ4] How does DES-LOC scale with increasing model size and longer training horizons?
\item[RQ5] How does DES-LOC perform when using a Nesterov outer optimizer?
\item[RQ6] Is DES-LOC compatible with non-Adam inner optimizers such as Muon?
\end{description}

%-----------------------------------------------------------------------------
\subsection{Experimental Setup}
\label{sec:experimental_setup}
%-----------------------------------------------------------------------------

\paragraph{Models and data.} We train a 135M-parameter GPT-style model (arch.\ in Table~\ref{tab:arch}) with sequence
length 2048. We distinguish worker batch size $B_w$ from global $B = \sum_{w=0}^{M-1} B_w$~\citep{sani2025photon}. A
2M token global batch is split across $M = 4$ workers sampling IID from SmolLM2~\citep{allal2025smollm2}:
70\% Fineweb-Edu~\citep{penedo2024fineweb}, 10\% Cosmopedia~\citep{benallal2024cosmopedia}, 10\%
Python-Edu, 5\% FineMath 4+, and 5\% Infi-WebMath 4+. The 135M model trains for
6.4B tokens ($2.4\times$ compute-optimal~\citep{hoffmann2022training}). For RQ4, we scale to a 1.7B model for
40B tokens ($2\times$ compute-optimal)~\citep{sardana2024beyond}. We show Non-IID data results in Fig.~\ref{fig:noniid}.

\paragraph{Optimizers.} We use Adam~\citep{kingma2015adam} (results in Section~\ref{sec:complementary}) and its variant
ADOPT~\citep{taniguchi2024adopt}, which modifies the update to guarantee convergence for any $\beta_2$.
For the 135M model, we grid-search $(\beta_1, \beta_2, \eta)$ under DDP; the 1.7B model uses hyperparameters
from~\citet{allal2025smollm2,taniguchi2024adopt}. Learning rates use the WSD schedule~\citep{hagele2024scaling,allal2025smollm2}. We favor ADOPT ($\beta_2 = 0.9999$) in high-$\beta$ regimes where Adam is unstable.
We also ablate the outer optimizer, comparing FedAvg with a Nesterov optimizer~\citep{reddi2021adaptive,douillard2023diloco,charles2025communication} on a 700M model trained on 40B tokens.

\paragraph{Baselines.} We compare DES-LOC with: (i) synchronous DDP; (ii) Local Adam/ADOPT; (iii)
FAVG+OPT (persistent states~\citep{sani2025photon,douillard2023diloco}); and (iv) FAVG$-$OPT
(reset states~\citep{sani2024future,iacob2025dept}). Persistent-state FedAvg is DES-LOC with
$K_u, K_v = \infty$, an upper bound on comms efficiency. DDP is an upper bound on ML performance.

\paragraph{Metrics.} We evaluate models by (i) perplexity and (ii) per-worker asymptotic communication cost
assuming a bandwidth-optimal Ring-AllReduce~\citep{sergeev2018horovod} algorithm scaling
linearly with model size. For the 1.7B model, we report standard in-context-learning (ICL) benchmarks~\citep{brown2020language}. We use a zero-shot setting for ICL tasks unless stated otherwise following~\citet{allal2025smollm2} and report the best performing communication-efficient method in blue with the
best-performing overall in bold. To fairly compare optimizer-state changes across decay rates, we
measure their \emph{relative} rates of change as $\|s_{t+K} - s_t\|_2 / \|s_t\|_2$. For convergence plots, we report
final-round means and standard deviations next to labels. We also provide wall-time clock results;
we use $4$ machines with one H100 for sub-1B models, and $4$ machines with $8$ H100s each for larger
scales. While the links between machines run at $100$Gb/s, we observed overheads limiting the
practical bandwidth to $60$--$70$ Gb/s. We report stepwise (see Section~\ref{sec:stepwise}) and timewise convergence.
We also provide an analysis on the wall-clock time vs bandwidth in Section~\ref{sec:wallclock_bandwidth}.

%=============================================================================
\section{Evaluation}
\label{sec:evaluation}
%=============================================================================

Our results show optimizer states change at different rates (Section~\ref{sec:rq1}), forming a clear synchronization hierarchy (Section~\ref{sec:rq2}). DES-LOC reduces communication $2\times$ vs.\ Local Adam (Section~\ref{sec:rq3})
while converging robustly with adding workers and scaling effectively to large models (Section~\ref{sec:rq4}).

%-----------------------------------------------------------------------------
\subsection{Higher $\beta$ Optimizer States Have Slower Empirical Rates of Change (RQ1)}
\label{sec:rq1}
%-----------------------------------------------------------------------------

Figure~\ref{fig:rates_of_change} shows that relative rates of change for the two momenta in Local ADOPT/Adam scale
with their decay rates under gradient clipping ($\rho = 1$). Supported by our theoretical discussions on
momenta half-lives (Section~\ref{sec:desloc}), the second momentum evolves substantially slower than the first at
high-$\beta_2$. For Local Adam, the second momentum remains slower even when $\beta_2 \approx \beta_1$, potentially
because gradient variance~\citep{kingma2015adam} evolves slower than the first momentum.

\paragraph{Takeaway:} As discussed in Sections~\ref{sec:desloc} and~\ref{sec:convergence}, when $\beta_1 \ll \beta_2$, the second momentum evolves slower
than the first, proportional to half-life ratio of the two $\frac{\tau_{0.5}(\beta_2)}{\tau_{0.5}(\beta_1)} = \frac{\ln(\beta_1)}{\ln(\beta_2)}$.

% Figure 2: PDF page 6 top
\begin{figure}[t]
\centering
% [Figure 2 graphic from original PDF]
\caption{Relative rates of change for first (left) and second (right) momenta across rounds using standard Local ADOPT ($K = 64$). For ADOPT ($\beta_2 = 0.9999$), increasing $\beta_1 \geq 0.99$ greatly slows the first-momentum rate of change. The second momentum evolves $\sim 100\times$ slower.}
\label{fig:rates_of_change}
\end{figure}

%-----------------------------------------------------------------------------
\subsection{Parameters Require Frequent Sync, Momenta Sync Proportional to $\beta$ (RQ2)}
\label{sec:rq2}
%-----------------------------------------------------------------------------

Figure~\ref{fig:sync_freq} evaluates the effect of independently varying synchronization periods ($K_x, K_u, K_v$) for
parameters and optimizer states. We consider two baseline periods ($K_b = 16, 256$), chosen based
on the fastest state's half-life ($\tau_{0.5}(0.95) \approx 13.5$). Frequent parameter synchronization ($K_x$) is
crucial for performance, while synchronizing momenta ($K_u, K_v$) significantly impacts training only
if their half-lives align with the base frequency $K_b$. Otherwise, synchronization frequency primarily
influences communication costs rather than model quality. Adam results can be seen in Section~\ref{sec:complementary}.

\paragraph{Takeaway:} Parameter synchronization frequency ($K_x$) strongly impacts performance, motivated
by the leading term in theoretical bounds (Section~\ref{sec:convergence}). Momentum synchronization periods matter
empirically only when chosen near their half-lives, consistent with Sections~\ref{sec:desloc} and~\ref{sec:convergence}.

% Figure 3: PDF page 7 top
\begin{figure}[t]
\centering
% [Figure 3 graphic from original PDF]
\caption{Model perplexity for DES-LOC (ADOPT, $\beta_1 = 0.95, \beta_2 = 0.9999$), varying synchronization periods independently (others fixed at $K_b$). Parameter synchronization (a) is critical, with degradation at higher periods. Second-momentum sync (b) minimally affects performance due to its large half-life. First-momentum sync improves perplexity when the freq matches its half-life; contrast (c) and (d). Parameters and second momentum behave similarly across frequencies (Section~\ref{sec:complementary})}
\label{fig:sync_freq}
\end{figure}

%-----------------------------------------------------------------------------
\subsection{DES-LOC Brings $2\times$ Communication Reductions over Local Adam (RQ3)}
\label{sec:rq3}
%-----------------------------------------------------------------------------

As shown in Figure~\ref{fig:comms}, DES-LOC halves communication versus Local Adam~\citep{cheng2025convergence}
with matching perplexity by syncing momenta less frequently ($K_u = 3K_x$, $K_v = 6K_x$),
exploiting the second-momentum's lower sensitivity to sync frequency (Fig.~\ref{fig:sync_freq}). This yields a $1.37\times$
speedup over DDP and $1.1\times$ over Local Adam at $K_x = 16$ on 4 H100s. At $K_x = 256$, the speedup
over DDP increases to $1.47\times$ and both DES-LOC and Local Adam saturate throughput.

As shown in Figure~\ref{fig:comms}, DES-LOC effectively initializes new workers and outperforms the straightforward approach of ad-hoc averaging from checkpoints---whose insufficiency is detailed in Section~\ref{sec:checkpoint}.

\paragraph{Takeaway:} DES-LOC achieves a $2\times$ communication reduction over Local Adam by leveraging two
insights: optimizer-state sync matters less than parameter sync, and slower-changing states (high $\beta_2$)
can sync less often. By eventually syncing all optimizer states, DES-LOC matches the robustness of
Local Adam with $K = \max(K_x, K_u, K_v)$ when adding new workers/responding to system failures.

% Figure 4: PDF page 7 bottom
\begin{figure}[t]
\centering
% [Figure 4 graphic from original PDF]
\caption{DES-LOC ($K_u, K_v = 3K_x, 6K_x$) reduces comms by $2\times$ over Local Adam while matching its performance and that of heuristic baselines at high (a) and low (b) frequencies (Section~\ref{sec:baselines}). We show robustness by doubling worker count at step $1536$ (c,d), where DES-LOC and Local Adam maintain stable perplexity/norms, outperforming heuristics and ad-hoc optimizer-state averaging.}
\label{fig:comms}
\end{figure}

%-----------------------------------------------------------------------------
\subsection{DES-LOC is Suitable for Large-Scale Training (RQ4)}
\label{sec:rq4}
%-----------------------------------------------------------------------------

Figure~\ref{fig:billion} shows that DES-LOC reliably scales to 1B models and very infrequent communication ($K_x = 256$). Evaluating the billion-scale models on the ICL tasks (Table~\ref{tab:icl}), DES-LOC is
competitive with all baselines while reducing communication versus Local Adam and DDP. The
heuristic baseline~\citep{sani2025photon} suffers training instabilities (Fig.~\ref{fig:billion}b) potentially impacting
downstream performance (Table~\ref{tab:icl}) and underscoring the advantage of DES-LOC's training stability.
We elaborate more on the settings in which we expect DES-LOC to improve stability in Section~\ref{sec:critical_batch}.

Our method's reduced communication costs result in a $\approx 2.2\times$ training speedup over DDP (Figure~\ref{fig:billion}).
As show by our benchmarks Table~\ref{tab:wallclock}, these time savings scale with model size and comms frequency.
At the $13$B scale with $K_x = 16$, DES-LOC would save 13 days over Local Adam and $73$ days
over DDP. The advantage over DDP widens for $K_x = 256$, where communication-efficient methods
maximize throughput. Table~\ref{tab:throughput} shows equivalent results for throughput.

% Figure 5: PDF page 8 top
\begin{figure}[t]
\centering
% [Figure 5 graphic from original PDF]
\caption{DES-LOC matches Local Adam perplexity (left) for billion-scale model training at half
the communication cost ($K_x = 256, K_u = 3K_x, K_v = 6K_x$), representing a $\mathbf{170\times}$ reduction
over DDP. Both DES-LOC and Local Adam converge to competitive perplexity at this scale.
FAVG+OPT achieves good performance (left) but suffers activation growth (right) and parameter-norm growth (Section~\ref{sec:complementary}), potentially due to noisy updates, raising concerns for extended training.}
\label{fig:billion}
\end{figure}

\begin{table}[t]
\centering
\caption{Our billion-scale model trained with DES-LOC matches or surpasses the ICL performance of models trained with Local Adam and FAVG+OPT, approaching DDP performance. FAVG+OPT
underperforms compared to its perplexity results from Fig.~\ref{fig:billion}a, indicating that the activation increases (Fig.~\ref{fig:billion}b) from the unstable training procedure may have damaged the model.}
\label{tab:icl}
\begin{tabular}{lcccc|c}
\toprule
& Arc Challenge & Arc Easy & PIQA & HellaSwag & Avg \\
\midrule
DES-LOC      & 31.8 & 59.0 & 70.7 & 44.9 & 51.6 \\
Local Adam   & 31.9 & 59.0 & 70.6 & 45.8 & 51.8 \\
FAVG+OPT     & 30.1 & 58.0 & 70.0 & 44.8 & 50.7 \\
\midrule
DDP          & 33.8 & 62.5 & 71.1 & 47.8 & 53.8 \\
\bottomrule
\end{tabular}
\end{table}

\begin{table}[t]
\centering
\caption{Wall-clock time (days) for 1B--13B models to reach $2\times$ compute-optimal tokens at high ($K = 16$) and low ($K = 256$) frequencies with a 2M token batch size. At high frequency ($K = 16$), DES-LOC outperforms Local ADAM by over 13 days on the 13B model and is within 3\% of FAVG+OPT. At low frequency ($K = 256$), it cuts the 13B's training time $>93$ days versus DDP.}
\label{tab:wallclock}
\begin{tabular}{l|cc|cc|cc}
\toprule
& \multicolumn{2}{c|}{1B Model} & \multicolumn{2}{c|}{7B Model} & \multicolumn{2}{c}{13B Model} \\
$K_x$ & 16 & 256 & 16 & 256 & 16 & 256 \\
\midrule
DDP (Baseline)   & $1.41 \pm 0.008$ & $1.41 \pm 0.008$ & $38.74 \pm 0.161$ & $38.74 \pm 0.161$ & $175.50 \pm 0.478$ & $175.50 \pm 0.478$ \\
\midrule
FAVG+OPT     & $0.80 \pm 0.007$ & $0.63 \pm 0.006$ & $28.52 \pm 0.095$ & $24.01 \pm 0.088$ & $100.21 \pm 0.544$ & $82.46 \pm 0.513$ \\
Local Adam   & $0.96 \pm 0.006$ & $0.64 \pm 0.006$ & $31.46 \pm 0.090$ & $24.18 \pm 0.087$ & $116.10 \pm 0.484$ & $83.34 \pm 0.509$ \\
DES-LOC ($K_u, K_v = 3K_x, 6K_x$) & $0.81 \pm 0.006$ & $0.63 \pm 0.006$ & $28.80 \pm 0.094$ & $24.06 \pm 0.088$ & $102.03 \pm 0.537$ & $82.68 \pm 0.512$ \\
\bottomrule
\end{tabular}
\end{table}

\paragraph{Takeaway:} DES-LOC enables efficient training of large-scale foundation models, especially at long
training horizons, with downstream ICL performance competitive with DDP. We recommend setting
$K_x$ for sufficient throughput based on bandwidth, then setting $K_u, K_v$ as constant multiples (e.g.,
$3\times$, $6\times$) or based on the half-life of their $\beta$ (see Section~\ref{sec:choosing_freq}).

%-----------------------------------------------------------------------------
\subsection{Nesterov as the Outer Optimizer (RQ5)}
\label{sec:rq5}
%-----------------------------------------------------------------------------

We ablate the outer optimizer for DES-LOC on a 700M parameter model, comparing averaging to a
Nesterov optimizer with momentum of 0.9, outer learning rate of 1.0 tuned following~\citet{charles2025communication}. The experiment ran on 4 H100s and used a medium-synchronization regime ($K_x = 32$, $K_u =
3K_x$, $K_v = 6K_x$) where models are initialised from 2048-step DDP checkpoints, following~\citet{charles2025communication}. While our convergence bound is not trivially applicable, our analysis of Eq.~(\ref{eq:psi}) suggests a higher momentum synchronization frequency ($p_u$) should permit a larger step size ($\eta_0 \propto 1/\sqrt{\psi}$).

As shown in Figure~\ref{fig:nesterov}, two key points emerge. First, more frequent synchronization ($K_x = 32$)
allows DES-LOC to come within $1\%$ of the final perplexity of DDP, performing much better than
in infrequent settings ($K_x = 256$). Second, using Nesterov as the outer optimizer improves
performance over averaging by $\approx 0.5\%$, with its performance w.r.t DDP being similar to the one
reported in~\citet{charles2025communication}, Table~4) for models at this scale. The Nesterov approach
preserves the practical benefits of DES-LOC, ensuring effective worker initialization and reducing
local optimization noise, which can help prevent issues like exploding activation norms.

% Figure 6: PDF page 9 top
\begin{figure}[t]
\centering
% [Figure 6 graphic from original PDF]
\caption{Ablation of the outer optimizer for DES-LOC on a $700$M parameter model in a medium-frequency communication setting ($K_x = 32$), showing (left) convergence in terms of time and (right) in terms of steps. In this regime, DES-LOC's final perplexity is within $1\%$ of the DDP baseline. Using a Nesterov outer optimizer provides a further improvement over averaging.}
\label{fig:nesterov}
\end{figure}

\paragraph{Takeaway:} Frequent synchronization ($K_x = 32$) allows DES-LOC to reach within 1\% of the perplexity of DDP. Furthermore, using a Nesterov outer optimizer improves performance over averaging,
while preserving practical benefits like effective worker initialization and reduced optimization noise.

We further investigate the benefits of synchronizing optimizer states under a Nesterov outer
optimizer relative to purely local-state methods (DiLoCo) at the 1B scale in Section~\ref{sec:nesterov_diloco}.

%-----------------------------------------------------------------------------
\subsection{Muon as the Inner Optimizer (RQ6)}
\label{sec:rq6}
%-----------------------------------------------------------------------------

To assess the versatility of DES-LOC beyond Adam-family optimizers, we integrate it with Muon~\citep{jordan2024muon}, a single-momentum optimizer that uses Newton-Schulz iterations for orthogonalization.
Since Muon preconditions the momentum term directly rather than tracking second-moment
estimates, the relevant synchronization periods reduce to parameters ($K_x$) and momentum ($K_u$). We
apply DES-LOC by synchronizing the Muon momentum buffer less frequently than the parameters,
setting $K_u = 3K_x$ with $K_x = 32$.

As shown in Figure~\ref{fig:muon}, DES-LOC matches the perplexity of the Local Muon baseline across
model scales ranging from 16M to $360$M parameters. By decoupling the synchronization frequencies,
DES-LOC communicates more than $1.5\times$ fewer bytes than the baseline, confirming that the Half-Life
Principle extends to optimizers beyond the Adam family. Full experimental details in Section~\ref{sec:muon_details}.

% Figure 7: PDF page 9 middle
\begin{figure}[t]
\centering
% [Figure 7 graphic from original PDF]
\caption{Training loss comparison between Local Muon ($K = 32$) and DES-LOC-Muon ($K_x = 32, K_u = 96$) across model scales ($16$M, $125$M, $360$M). DES-LOC matches the Local Muon baseline across all scales while communicating more than $\mathbf{1.5\times}$ fewer bytes. See Section~\ref{sec:muon_details} for full experimental details.}
\label{fig:muon}
\end{figure}

\paragraph{Takeaway:} DES-LOC is compatible with single-momentum optimizers such as Muon that rely
on Newton-Schulz preconditioning. By reducing the synchronization frequency of the momentum
buffer relative to parameters, DES-LOC maintains solution quality while significantly lowering
communication volume, demonstrating that our approach is inner-optimizer-agnostic.

%=============================================================================
\section{Related Work}
\label{sec:related}
%=============================================================================

In synchronous data-parallel training, workers exchange full gradients or parameters \textit{every} iteration,
incurring communication costs linear in model size using Ring-AllReduce~\citep{sergeev2018horovod}.
When hardware is weakly connected or widely distributed, communication significantly
slows wall-clock training time~\citep{sani2025photon} as workers need to wait for synchronization to
finish. Federated Averaging (FedAvg)~\citep{mcmahan2017communication} and Local SGD~\citep{stich2019local}
reduce communication by performing $K$ local optimization steps before averaging parameters,
decreasing communication rounds by a factor of $K$. Ad-hoc extensions to adaptive optimizers either
keep optimizer states local~\citep{douillard2023diloco,charles2025scaling,liu2024asynchronous} or reset them
after each sync~\citep{sani2024future,sani2025photon}, both lacking robust convergence guarantees.

Adam~\citep{kingma2015adam} is popular for pre-training as it scales to larger batches than SGD~\citep{kunstner2023noise,dubey2024llama}. It uses moving averages of gradients and their squares, however,
its convergence is not guaranteed as it requires $\beta_1 < \sqrt{\beta_2} < 1$, with large, problem-specific $\beta_2$~\citep{reddi2018convergence,zhang2022adam}. Other optimizers also track gradient moments~\citep{sutskever2013importance,chen2023symbolic,you2020large,taniguchi2024improved}. \textit{Local Adam}~\citep{cheng2025convergence} reduces communication with local steps but requires syncing optimizer states, which triples the
communication cost relative to Local SGD/DDP, as sync costs scale with the number of states. For
further related work, including compression/sparsification and structured updates, check Section~\ref{sec:extended_related_work}.

%=============================================================================
\section{Conclusion}
\label{sec:conclusion}
%=============================================================================

DES-LOC reconciles communication efficiency with rigorous convergence guarantees in distributed
adaptive optimization. By extending theory to the independent synchronization of Adam and SGDM
optimizer states, we empirically demonstrate convergence alongside $\mathbf{170\times}$ and $\mathbf{2\times}$ communication
reductions over DDP and prior state-of-the-art methods at billion-scale LLM training, even in environments prone to system failures. Our findings yield clear guidelines: i) \textbf{frequently} synchronize
parameters, and ii) synchronize optimizer states \textbf{less often}, proportional to their half-lives. These
insights open avenues for future research, including layer-wise synchronization, adaptive frequencies,
compressed updates, as well as emerging applications, such as worldwide cross-data center training
and collaborative training. As training workloads scale, we envision DES-LOC becoming the standard
for efficient, resilient foundation-model training in data centers and distributed environments.

\section*{Acknowledgments}
All costs for the computational resources used for this work were funded by Flower Labs, and the
research conducted by a team of researchers from Flower Labs, The University of Cambridge, The
Institute of Science and Technology of Austria, The University of Warwick, and Mohamed bin
Zayed University of Artificial Intelligence. Support for university-based researchers came from
a variety of sources, but in particular, the following funding organizations are acknowledged: the
European Research Council (REDIAL), the Royal Academy of Engineering (DANTE), the Ministry of
Education of Romania through the Credit and Scholarship Agency, and the European Union's Horizon
2020 research and innovation programme under the Marie Sk\l{}odowska-Curie grant agreement No
101034413.

\bibliography{references}
\bibliographystyle{neurips_2026}




\appendix

%=============================================================================
\section{Experimental Details and Optimizer Hyperparameter Sweeps}
\label{sec:appendix_a}
%=============================================================================

Here we provide additional experimental details complementing those in Section~4.1, including: a) model architecture details and hyperparameters independent of optimizer choice (Section~\ref{sec:arch_details}), b) our hyperparameter sweep procedure to select optimizer-specific settings (Section~\ref{sec:sweep}), and c) the optimal hyperparameters with those used in Section~5 highlighted in bold.

\subsection{Architecture Details and Hyperparameters}
\label{sec:arch_details}

\begin{table}[h]
\centering
\caption{Model architecture and training parameters. We denote the number of transformer blocks by \#Blocks, number of attention heads by \#Heads, embedding dimension by $d_{\text{model}}$, vocabulary size by $|V|$, and feedforward-layer expansion by Exp.\ Ratio. All models use positional embeddings~\citep{su2024roformer}, the silu activation function, and norm-based gradient clipping with clip-bound $\rho$. Global batch size (summed across all workers) is $|B_G|$, and sequence length is standard for models at these scales. For model initialization we use $\sigma = 1/\sqrt{d_{\text{model}}}$. The total number of steps is denoted by $T$.}
\label{tab:arch}
\begin{tabular}{lccccccccccc}
\toprule
Model Size & Blocks & $d_{\text{model}}$ & $|V|$ & \#Heads & Exp.\ Ratio & ROPE $\theta$ & ACT & Init $\sigma$ & $\rho$ & Seq Len & $|B_G|$ & $T$ \\
\midrule
135M & 30 & 576  & 50K & 9  & 4 & 10000 & silu & 0.04 & 1.0 & 2048 & 1024 & 1536, 3072 \\
720M & 12 & 2048 & 50K & 16 & 4 & 10000 & SiLU & 0.02 & 1.0 & 2048 & 512  & 38912 \\
1.7B & 24 & 2048 & 50K & 16 & 4 & 10000 & silu & 0.02 & 1.0 & 2048 & 1024 & 20480 \\
\bottomrule
\end{tabular}
\end{table}

Table~\ref{tab:arch} summarizes the architectural details of our models, following established practices for large
language models at their respective scales. Unless otherwise stated, we adopt the hyperparameters
recommended by~\citet{allal2025smollm2} for both the 135M and the 1.7B models. We operate at a
batch size of 2M tokens, which is very large for the 135M model at the length of training we
perform~\citep{zhang2025critical} and industry-standard for the 1.7B model~\citep{touvron2023llama2}, we
chose to operate at large batch sizes because adaptive optimizers provide benefits primarily in large-batch training regimes~\citep{kunstner2023noise}. Moreover, we intend \texttt{DES-LOC} for use in cross data-center scenarios, where effectively utilizing available accelerators naturally demands large batch sizes and/or model scales. For both model sizes, we train for approximately $2\times$ the compute-optimal token budget~\citep{hoffmann2022training}, placing our evaluations within the context of extended-duration foundation model training~\citep{allal2025smollm2}. Our chosen token budget is conservative due to resource constraints; for comparison, \citet{allal2025smollm2} used 11 trillion tokens which is over $4000\times$ compute-optimal for the 135M model, and $300\times$ for the 1.7B.

We select warmup and decay schedules following recommendations from~\citet{zhang2025critical,hagele2024scaling,allal2025smollm2}. For the 135M model, the warmup period is set to $T_{\text{WARM}} = 512$
steps, corresponding to roughly 40\% of the compute-optimal training tokens recommended by~\citet{zhang2025critical}. For the 1.7B model, we use the recommended $T_{\text{WARM}} = 2048$ steps from~\citet{allal2025smollm2}, roughly 10\% of total training. The stable-decay period uses a $1 - \text{SQRT}$ schedule over
the final $T_{\text{DECAY}} = 10\% \times T$ steps~\citep{hagele2024scaling}. For shorter runs, such as $T = 1536$ during heterogeneous-data evaluations, we keep the warmup fixed and proportionally scale the decay to ensure well-conditioned parameter updates during the stable learning rate period. The seeds we use for data sampling and for controlling the training algorithms and model are provided in the code accompanying the appendix.

\subsection{Optimizer Parameters Sweeping Procedure}
\label{sec:sweep}

As detailed in Section~\ref{sec:desloc} and verified empirically in Section~\ref{sec:rq2}, the choice of decay rates $\beta_1, \beta_2$
strongly influences the effective synchronization frequencies achievable by both DES-LOC and
Local Adam. This relationship arises directly from the half-life of optimizer states, given by
$\tau_{0.5} = \frac{\ln(0.5)}{\ln(\beta)}$.

For Adam, prior studies such as \citet{wortsman2023stable} have demonstrated a critical interplay between the learning rate ($\eta$), batch size, and the second-momentum decay $\beta_2$. Specifically, increasing either the learning rate or batch size typically demands a lower $\beta_2$ to maintain training stability and avoid loss spikes. Conversely, higher $\beta_2$ values constrain the learning rate and batch size. Such dynamics have also been recently observed between the learning rate and the first-momentum decay $\beta_1$ in~\citet{pagliardini2025ademamix}. Given that all our experiments use a fixed large batch size of roughly 2 million tokens (appropriate for billion-scale training), we systematically tune the learning rate $\eta$ in response to changes in $\beta_1, \beta_2$. We try values of $\beta_1, \beta_2$ based on previous works~\citep{zhang2025critical} and follow the theoretical convergence requirement of~\citet{zhang2022adam} setting $\beta_1 \leq \sqrt{\beta_2}$.

Due to computational constraints, we cannot jointly optimize synchronization periods, data distributions, and decay parameters, and instead adopt a structured two-stage tuning approach:
\begin{enumerate}
    \item \textbf{Stage 1: Tuning $\eta$ for DDP.} Starting from the recommended baseline learning rate ($\eta_0$)
    from~\citet{allal2025smollm2}, we conduct a grid search as outlined by~\citet{charles2025communication}:
    $\{\ldots, \sqrt{2}^{-2}\eta_0, \sqrt{2}^{-1}\eta_0, \eta_0, \sqrt{2}\eta_0, \sqrt{2}^2\eta_0, \ldots\}$
    We expand this search until perplexity stops improving, identifying an optimal learning rate $\eta^*_{\text{DDP}}$ for each $(\beta_1, \beta_2)$ configuration.

    \item \textbf{Stage 2: Tuning $\eta$ for Local Adam.} We then repeat this procedure for Local Adam,
    using $\eta^*_{\text{DDP}}$ as the new baseline. To balance generalizability and computational cost, we set
    the synchronization period to an intermediate value of $K = 64$, between high-frequency ($K = 16$) and low-frequency ($K = 256$) scenarios.
\end{enumerate}

Additionally, following~\citet{zhang2025critical}, we omit weight decay (set to zero) to simplify the hyperparameter tuning process, as it directly affects only model parameters, not optimizer states.

For experiments using Nesterov, we follow the hyperparameter sweeping procedure of~\citet{charles2025communication}, starting with a server learning rate of $1.0$ and a momentum of $0.9$ and only lowering it if it fails to converge.

\subsubsection{Optimizers' Hyperparameter Configurations}
\label{sec:hp_configs}

\begin{table}[h]
\centering
\caption{Optimal learning rates $\eta^*$ for $\beta_1, \beta_2$ configurations of ADOPT/Adam. The hyperparameter sweep procedure (see Section~\ref{sec:sweep}) involves incrementally adjusting the learning rate by factors of $\sqrt{2}$ around the initial value from~\citet{allal2025smollm2} until performance stops improving.}
\label{tab:lr_sweep}
\begin{tabular}{llll}
\toprule
Optimizer & $\beta_1$ & $\beta_2$ & $\eta^*$ \\
\midrule
\multirow{4}{*}{ADOPT} & 0.9   & 0.9999 & 0.0021 \\
                        & \textbf{0.95}  & \textbf{0.9999} & \textbf{0.0021} \\
                        & 0.99  & 0.9999 & 0.0014 \\
                        & 0.995 & 0.9999 & 0.0007 \\
\midrule
\multirow{5}{*}{Adam}  & 0.9   & 0.95   & 0.0042 \\
                        & \textbf{0.95}  & \textbf{0.95}   & \textbf{0.003}  \\
                        & 0.9   & 0.99   & 0.003  \\
                        & 0.95  & 0.99   & 0.003  \\
                        & 0.99  & 0.99   & 0.0021 \\
\bottomrule
\end{tabular}
\end{table}

Our hyperparameter sweep (Table~\ref{tab:lr_sweep}) indicates that the optimal learning rate $\eta^*$ under the warmup-stable-decay scheduler~\citep{hagele2024scaling} strongly depends on both optimizer type and the chosen $\beta_1, \beta_2$ values. For Adam, optimal learning rates and second-momentum decay ($\beta_2$) align closely with recommendations from~\citet{allal2025smollm2}, though a slightly higher first-momentum decay ($\beta_1$) consistently performs better, in agreement with prior findings~\citep{zhang2025critical}. For ADOPT (default $\beta_2$), we observe a lower optimal learning rate compared to Adam, but similar best-performing $\beta_1$ values. We also find that the optimal learning rate does not differ between DDP and Local Adam for given $\beta_1, \beta_2$ when $K = 64$ and using a $\sqrt{2}$ sweep, higher learning rates either do not provide a benefit or diverge while lower learning rates are only necessary when pushing $K$ far closer to the complete training duration.

We find that increasing $\beta_1$ for ADOPT, and $\beta_1, \beta_2$ for Adam, leads to rapid performance degradation, particularly at or above $0.99$. Since the half-life at $\beta = 0.99$ ($\tau_{0.5} \approx 69$) is not sufficiently longer than at $\beta = 0.95$ ($\tau_{0.5} \approx 13.5$) to justify the observed performance drop, we select $\beta_1 = 0.95$ for all experiments, along with the default $\beta_2$ for ADOPT and $\beta_2 = 0.95$ for Adam.

\paragraph{Takeaway.} Increasing an optimizer state's $\beta$ significantly affects performance. Since linear increases in $\beta$ cause only logarithmic changes in half-life $\tau_{0.5}$, raising $\beta$ beyond the optimal value degrades performance without substantially improving the achievable synchronization frequency (Section~\ref{sec:rq2}).

%=============================================================================
\section{Complementary Results to Sections 2.1 and 5}
\label{sec:complementary}
%=============================================================================

We now provide additional results supplementing those presented in the main text. Specifically:

\begin{enumerate}
\item \textbf{Section B.1} complements Section~2.1 by including results on the heterogeneous data distribution described in Section~4.1. This highlights \texttt{DES-LOC}'s robustness under imperfect sampling or strongly Non-IID federated scenarios (see Kairouz et al., 2021, Sec 3.1).
\item \textbf{Section B.2.1} complements Fig.~3 by showing the separate impact of varying synchronization frequencies for parameters and the second momentum when the base frequency is $K_b=16$. It supports our claim that parameters and second momentum exhibit similar behavior across different synchronization regimes, unlike the first momentum.
\item \textbf{Section B.2.2} extends Fig.~3 by evaluating \texttt{DES-LOC-Adam}. We confirm that the parameter synchronization frequency is the most important, as predicted by our theory. In contrast, the momenta sync frequency is far less impactful, especially for low parameter sync frequencies.
\item \textbf{Section B.3.2} complements Fig.~4 by showing \texttt{DES-LOC-ADOPT}'s perplexity against baseline methods on heterogeneous data (as defined in Section~4.1). This validates our claim from Contribution~2 regarding \texttt{DES-LOC}'s effectiveness on heterogeneous datasets.
\item \textbf{Section B.3.3} presents an ablation study examining alternative low-communication configurations of \texttt{DES-LOC}, justifying our choice of $K_u=3K_x, K_v=6K_x$ used in Fig.~4.
\item \textbf{Section B.3.4} repeats the baseline comparison from Fig.~4 for \texttt{DES-LOC-Adam}, demonstrating that \texttt{DES-LOC} achieves similar communication reductions and performance when using Adam instead of ADOPT.
\item \textbf{Section B.4} provides additional metrics illustrating training instabilities for the \texttt{FAVG+OPT} baseline, including rapidly growing parameter norms, supporting observations in Fig.~5.b.
\item \textbf{Section B.5} presents the full 1B-scale comparison of \texttt{DES-LOC} Nesterov against DiLoCo, demonstrating the benefits of periodic optimizer-state synchronization under a Nesterov outer optimizer.
\item \textbf{Section B.6} benchmarks \texttt{DES-LOC} under extremely low bandwidth ($10$ Gbit/s), showing $9.42\times$ wall-clock speedups over DDP.
\item \textbf{Section B.7} provides full experimental details for the Muon experiments presented in Section~5.6.
\item \textbf{Section B.8} evaluates \texttt{DES-LOC} on the Flux vision model, demonstrating generality beyond LLMs.
\item \textbf{Section B.9} measures training throughput at the $7$B parameter scale on B200 GPUs.
\end{enumerate}

%-----------------------------------------------------------------------------
\subsection{Toy Problem on Non-IID Data (See Section~2.1)}
\label{sec:b1_toy_noniid}
%-----------------------------------------------------------------------------

\textbf{Toy Example Non-IID:} Fig.~8 simulates the scenario from Section~3, where each worker $m$ optimizes a distinct loss $f_m$ on heterogeneous data. Both \texttt{DES-LOC} and Local Adam show more stable convergence and get closer to the optimum than heuristic baselines.

% Figure 8: Toy problem Non-IID (PDF page 21 top)
\begin{figure}[ht]
\centering
% [placeholder for toy problem Non-IID figure]
\caption{We present a toy problem in a Non-IID setting, where \texttt{DES-LOC} (with synchronization periods $K_x=192, K_u=192, K_v=692$) and Local Adam (with $K=K_x$) converge to a superior solution compared to methods that keep optimizer states local \citep{Douillard2023,Sani2025} or periodically reset them \citep{Sani2024,Iacob2025}. Like the IID scenario, resetting optimizer states prevents convergence due to repeated oscillations caused by reinitializations. Additionally, as seen in panel~(a) between rounds $15$ and $40$, methods keeping optimizer states local suffer from larger oscillations further away from the optimum. The function optimized is $f(x_1,x_2)=(1-x_1)^2+100(x_2-x_1^2)^2$, and we simulate $M=256$ workers, each adding Gaussian noise with worker-specific standard deviation $\sigma^m \sim \mathcal{N}(0,3)$.}
\label{fig:8}
\end{figure}

%-----------------------------------------------------------------------------
\subsection{RQ2: Independent Sync Frequencies}
\label{sec:b2_sync_freq}
%-----------------------------------------------------------------------------

This section provides supplementary results for \textbf{RQ2}, complementing Section~5.2. Section~B.2.1 shows that perplexity has similar sensitivity to the first and second momentum synchronization frequencies at both high and low base synchronization frequencies. Additionally, Section~B.2.2 repeats the comparison from Fig.~3 for \texttt{DES-LOC-Adam}, revealing similar trends regarding the importance of the parameters, with a reduced importance for the momenta due to lower $\beta_2$.

\subsubsection{Parameter and Second Momentum At $K_b=16$ (See Fig.~3.a, Fig.~3.b)}
\label{sec:b21_param_second}

Figure~9 examines the effects of independently varying synchronization periods ($K_x, K_v$) for parameters and second momentum under \texttt{DES-LOC-ADOPT} in the high-frequency regime ($K_b=16$), chosen based on the first momentum's half-life ($\tau_{0.5}\approx 13.5$). Similar to the low-frequency results in Fig.~3.a, parameter synchronization frequency ($K_x$) strongly influences perplexity, while the second momentum ($K_v$) has minimal impact due to its long half-life. This contrasts with the first momentum, whose half-life closely matches the high-frequency period.

% Figure 9: DES-LOC ADOPT perplexity at Kb=16 (PDF page 21 bottom)
\begin{figure}[ht]
\centering
% [placeholder for Fig 9]
\caption{Model perplexity for \texttt{DES-LOC} (ADOPT, $\beta_1=0.95, \beta_2=0.9999$), independently varying synchronization periods at a high baseline frequency ($K_b=16$). Similar to Fig.~3, parameter synchronization~(a) is critical, with performance sharply degrading at higher periods, while second-momentum synchronization~(b) has minimal impact due to its large half-life ($\tau_{0.5}(\beta_2)\gg K_b$).}
\label{fig:9}
\end{figure}

\textbf{Takeaway:} In high-frequency synchronization regimes, the importance of parameters and the second momentum remains similar to the low-frequency regime shown in Section~5.2.

\subsubsection{Adam Results (See Fig.~3)}
\label{sec:b22_adam}

Fig.~10 show the rate of change results for Adam momenta at various $\beta$ values.

Figure~11 provides complementary results to Figs.~3 and~9 using \texttt{DES-LOC-Adam} with $\beta_1=\beta_2=0.95$. Unlike ADOPT, the relatively low $\beta$ result in both the first and second momentum quickly adapting to the local gradients, reducing the impact of their sync frequency.

% Figure 10: Rate of change for Adam momenta (PDF page 22)
\begin{figure}[ht]
\centering
% [placeholder for Fig 10]
\caption{Relative rates of change for first and second momenta across rounds using standard Local Adam ($K=64$). Increasing $\beta_1$ substantially reduces the rate of change of the first momentum, while increasing either $\beta_1, \beta_2$ decreases the rate of change of the second.}
\label{fig:10}
\end{figure}

% Figure 11: DES-LOC-Adam perplexity varying sync periods (PDF page 22)
\begin{figure}[ht]
\centering
% [placeholder for Fig 11]
\caption{Model perplexity for \texttt{DES-LOC-Adam} ($\beta_1=\beta_2=0.95$) when independently varying sync periods ($K_x, K_u, K_v$) while fixing others at baseline $K_b$. Parameter synchronization~(a,b) influences performance in both high ($K_b=16$) and low ($K_b=256$) frequency regimes. Momenta synchronization minimally impacts perplexity due to both states' high adaptivity (low $\beta$), with potentially minor effects during the early stages of training in high-frequency regimes~(c,e).}
\label{fig:11}
\end{figure}

\textbf{Takeaway:} For \texttt{DES-LOC-Adam}, parameter synchronization remains critical, consistent with theory. However, due to reduced $\beta_2$, momenta synchronization is less impactful since both the numerator and denominator of Adam updates are driven by local worker gradients after a few initial steps.

%-----------------------------------------------------------------------------
\subsection{RQ3: Communication Reduction And Baseline Comparisons}
\label{sec:b3_comm_reduction}
%-----------------------------------------------------------------------------

This section provides supplementary results for \textbf{RQ3}, complementing Section~5.3. Section~B.3.2 shows the perplexity of different configurations providing a $2\times$ communication reduction over Local Adam. Additionally, Section~B.3.4 repeats the comparison against baselines from Section~5.3 for \texttt{DES-LOC-Adam}, showing similar communication reductions relative to Local Adam.

\subsubsection{Stepwise plots for Baseline Comparison}
\label{sec:b31_stepwise}

Figures~12 and~13 show stepwise plots for wall-clock results in the main text, they are the counterparts to Figs.~4 and~5.

% Figure 12 and 13: Stepwise plots (PDF page 23)
% [placeholder for Figs 12-13 -- stepwise counterparts to Figs 4,5]

\subsubsection{DES-LOC on Heterogeneous Data (See Contribution~2)}
\label{sec:b32_heterogeneous}

Figure~14 evaluates the robustness of \texttt{DES-LOC} against baselines under heterogeneous (Non-IID) data distributions as described in Section~4.1. We set synchronization periods to $K_x=K$, $K_u=3K_x$, and $K_v=6K_x$ to achieve a targeted $2\times$ communication reduction over Local Adam.

% Figure 14: Non-IID perplexity comparison (PDF page 24 top)
\begin{figure}[ht]
\centering
% [placeholder for Fig 14]
\caption{Comparison of perplexity under Non-IID conditions for \texttt{DES-LOC}, Local Adam ($K_x=K_u=K_v$), and heuristic baselines (defined in Section~4.1) at high~(a) and low~(b) synchronization frequencies. Due to higher cross-worker variance caused by heterogeneous data, parameters require slightly more frequent synchronization in the low-frequency regime ($K_x=128<256$). Experiments are limited to $T=1536$ steps ($\sim$compute-optimal) for computational feasibility.}
\label{fig:14}
\end{figure}

\subsubsection{DES-LOC Low Communication Configurations Ablation (See Fig.~4)}
\label{sec:b33_ablation}

Figure~15 explores alternative synchronization configurations enabling \texttt{DES-LOC} to achieve improved communication efficiency over Local Adam. Motivated by theoretical insights (Sections~2 and~3) and empirical evidence (Sections~5.1 and~5.2), we only consider settings where parameter synchronization is most frequent ($K_x \leq \min(K_u, K_v)$). This constraint follows from experiments in Section~5.2, which show that infrequent parameter synchronization significantly degrades perplexity, while momentum synchronization frequency has a smaller impact. For a fixed $2\times$ communication reduction over Local Adam, our findings confirm that synchronizing the first momentum more frequently than the second aligns with their respective half-lives and maintains performance close to Local Adam.

% Figure 15: DES-LOC 2x communication configs ablation (PDF page 24 bottom)
\begin{figure}[ht]
\centering
% [placeholder for Fig 15]
\caption{Configurations of \texttt{DES-LOC} targeting $2\times$ lower communication than Local Adam ($K_x=K_u=K_v$), setting $K_u, K_v$ as multiples of $K_x$. In both high~(a) and low-frequency~(b) regimes, performance depends on how communication is split between momenta for $\beta_1 \ll \beta_2$. Syncing the first momentum less often ($K_u=6K_x, K_v=3K_x$) degrades performance, wasting communication on the slow second momentum. Conversely, syncing it frequently ($K_u=3K_x, K_v=6K_x$) yields performance comparable to Local Adam. Setting $K_u=K_v=4K_x$ produces intermediate results.}
\label{fig:15}
\end{figure}

\subsubsection{Adam Results (See Fig.~4)}
\label{sec:b34_adam_fig4}

We now present results for \texttt{DES-LOC-Adam} with $\beta_1=\beta_2=0.95$. \texttt{DES-LOC-Adam} achieves similar communication reductions over Local Adam and DDP as ADOPT. However, due to the lower $\beta_2$, the second-momentum half-life ($\tau_{0.5}(0.95)\approx 13.5$) is significantly shorter than for ADOPT ($\tau_{0.5}(0.9999)\approx 6931$). Figure~16 shows that with both momenta evolving at similar rates, the benefit of selecting $K_u < K_v$ diminishes. For consistency and due to meaningful empirical differences in rates of change (Section~5.1), we keep $K_u=3\times K_x$ and $K_v=6\times K_x$ in subsequent comparisons.

Figure~17 shows \texttt{DES-LOC-Adam} achieves a $2\times$ communication reduction over the prior state-of-the-art Local Adam \citep{Cheng2025} without significant perplexity degradation. Due to the much faster evolution of the optimizer states using Adam compared to ADOPT, local worker gradients drive the optimization reducing the benefit of allocating more of the communication budget to the first momentum.

% Figure 16: DES-LOC-Adam 2x communication configs (PDF page 25)
\begin{figure}[ht]
\centering
% [placeholder for Fig 16]
\caption{Configurations of \texttt{DES-LOC} targeting $2\times$ lower communication than Local Adam ($K_x=K_u=K_v$), using Adam ($\beta_1=\beta_2=0.95$). In contrast to \texttt{DES-LOC-ADOPT} (where $\beta_1 \ll \beta_2$ yields an advantage for $K_u < K_v$ as shown in Fig.~15), the similar half-lives in Adam make perplexity insensitive to how communication is split between momenta for high~(a) and low-frequencies~(b).}
\label{fig:16}
\end{figure}

% Figure 17: DES-LOC-Adam baseline comparison (PDF page 25)
\begin{figure}[ht]
\centering
% [placeholder for Fig 17]
\caption{\texttt{DES-LOC-Adam} achieves a $2\times$ communication reduction over the prior state-of-the-art Local Adam without significant perplexity degradation.}
\label{fig:17}
\end{figure}

\textbf{Takeaway:} \texttt{DES-LOC-Adam} achieves a similar $2\times$ communication reduction over Local Adam as \texttt{DES-LOC-ADOPT} by exploiting the reduced importance of optimizer-state synchronization relative to parameters. However, due to the smaller $\beta_2$ in Adam, there is limited benefit from assigning different synchronization frequencies to the first and second momenta compared to ADOPT.

%-----------------------------------------------------------------------------
\subsection{RQ4: Additional Metrics and Training Instabilities of FAVG+OPT (See Fig.~5.b)}
\label{sec:b4_favgopt}
%-----------------------------------------------------------------------------

Figure~18 complements Fig.~5.b by showing parameter and update norms for \texttt{DES-LOC} and baseline methods when training billion-scale models. Both \texttt{DES-LOC} and Local Adam regularize updates by synchronizing optimizer states, effectively reducing update norms due to averaging across workers (triangle inequality). In contrast, the heuristic baseline \citep{Sani2025} experiences large updates, leading to uncontrolled parameter growth, increased activations (Fig.~5.b), and degraded performance on downstream ICL tasks (Table~1) relative to its perplexity (Fig.~5.a).

% Figure 18: Update and parameter norms comparison (PDF page 26)
\begin{figure}[ht]
\centering
% [placeholder for Fig 18]
\caption{Comparison of update~(a) and parameter norms~(b) for billion-scale models trained with \texttt{DES-LOC} ($K_x=256, K_u=768, K_v=1536$), Local Adam ($K=256$), DDP, and Federated Averaging with persistent optimizer states (\texttt{FAVG+OPT}). Frequent synchronization in Local Adam and DDP consistently reduces update and parameter norms. Similarly, \texttt{DES-LOC} achieves comparable reductions at intervals corresponding to multiples of $\texttt{lcm}(K_x, K_u, K_v)$, with smaller intermediate drops. Conversely, \texttt{FAVG+OPT}, which does not synchronize optimizer states, experiences persistently larger and noisier updates, becoming vulnerable to spikes (notably before step $5000$). This leads to uncontrolled parameter growth~(b).}
\label{fig:18}
\end{figure}

% Table 5: Throughput comparison (PDF page 26, nougat inline)
\begin{table}[ht]
\centering
\caption{Throughput (batches/sec) for 1B--13B models to reach $2\times$ compute-optimal tokens at high ($K=16$) and low ($K=256$) frequencies with a 2M token batch size. All local methods achieve significant throughput gains over the DDP baseline. At high frequency ($K=16$), \texttt{DES-LOC} boosts throughput by over $1.7\times$ on the $13$B model. At low frequency ($K=256$), this advantage grows to over $2.1\times$ versus DDP.}
\label{tab:5}
\begin{tabular}{l c c c c c c}
\hline \hline
 & \multicolumn{2}{c}{\textbf{1B Model}} & \multicolumn{2}{c}{\textbf{7B Model}} & \multicolumn{2}{c}{\textbf{13B Model}} \\
\cline{2-7}
\textbf{Method} & $K_x\!=\!16$ & $K_x\!=\!256$ & $K_x\!=\!16$ & $K_x\!=\!256$ & $K_x\!=\!16$ & $K_x\!=\!256$ \\
\hline
DDP (Baseline) & $171.9 \pm 0.94$ & $171.9 \pm 0.9$ & $25.1 \pm 0.10$ & $25.1 \pm 0.10$ & $11.1 \pm 0.03$ & $11.11 \pm 0.03$ \\
\hline
FAVG+OPT & $304.9 \pm 2.51$ & $385.4 \pm 3.7$ & $34.0 \pm 0.11$ & $40.4 \pm 0.15$ & $19.4 \pm 0.11$ & $23.5 \pm 0.15$ \\
Local Adam & $253.1 \pm 1.59$ & $380.1 \pm 3.6$ & $30.9 \pm 0.09$ & $40.2 \pm 0.15$ & $16.7 \pm 0.07$ & $23.3 \pm 0.14$ \\
DES-LOC ($K_u,K_v\!=\!3K_x,6K_x$) & $299.2 \pm 2.39$ & $384.1 \pm 3.7$ & $33.7 \pm 0.11$ & $40.4 \pm 0.15$ & $19.0 \pm 0.10$ & $23.5 \pm 0.15$ \\
\hline \hline
\end{tabular}
\end{table}

\textbf{Takeaway:} Unlike heuristic methods, which maintain purely local optimizer states leading to unstable, noisy updates, \texttt{DES-LOC} provides stable regularization similar to Local Adam and DDP by periodically synchronizing parameters and momenta, reducing training instabilities.

%=============================================================================
\subsection{DES-LOC Nesterov vs.\ DiLoCo at 1B Scale}
\label{sec:b5_nesterov_diloco}
%=============================================================================

Having shown that a Nesterov outer optimizer improves \texttt{DES-LOC} (Section~5.5), we ask whether synchronizing optimizer states still helps relative to the local-state ($K_u=K_v=\infty$) Nesterov method DiLoCo. \citet{charles2025convergence} has shown that at \textbf{large} scale ($>1B$ parameters) DiLoCo can match or outperform DDP with Adam. We adopt their outer hyper-params and train a 1B model with the same experimental design as Section~5.5 using $4\times 8$ H100s for $40{,}960$ steps ($\approx 4\times$ compute-optimal) with inner Adam, comparing Adam-DDP, Local Adam, DiLoCo, and \texttt{DES-LOC} Nesterov, with $K_x=32$ for Local Adam and DiLoCo and $K_u=4K_x, K_v=8K_x$ for \texttt{DES-LOC} Nesterov.

\begin{figure}[t]
\centering
% Figure 19: 4-panel comparison (PDF page 27 top)
\caption{Comparison of Adam-DDP, Local Adam, DiLoCo, and \texttt{DES-LOC} Nesterov on a 1B-parameter model trained for $40{,}960$ steps with an AdamW inner optimizer. Top left: train perplexity vs steps. Top right: worker gradient norms. Bottom left: train perplexity vs time, bottom right: whole-model output activation norms. Shaded regions show std across workers. \texttt{DES-LOC} Nesterov outperforms Adam, it also outperforms DiLoCo at the cost of more communication. Error bars show variance across workers, accounting for compounding local drift.}
\label{fig:nesterov_diloco_1b}
\end{figure}

Figure~\ref{fig:nesterov_diloco_1b} (top left) shows that Nesterov-based methods outperform AdamW: DiLoCo achieves $7.63\pm 0.20$ validation perplexity, improving over Adam-DDP by $\approx 2\%$, while \texttt{DES-LOC} Nesterov reaches $7.56\pm 0.20$, a $\approx 0.9\%$ gain over DiLoCo; both outperform Local Adam ($8.13\pm 0.23$). Note that this comparison pits Nesterov-based local updates against DDP with standard Adam; as \citet{charles2025convergence} note, once DDP also uses Nesterov, this gap can shrink or reverse depending on model size and worker count. These results show that synchronizing optimizer states preserves the benefits of Nesterov while retaining the advantages of state averaging.

We analyze the interaction between the optimizer states and the outer optimizer by measuring the gradient norms and activation statistics. In Figure~\ref{fig:nesterov_diloco_1b} (top and bottom right), for both DiLoCo and \texttt{DES-LOC} Nesterov, gradient norms drop rapidly relative to Adam-DDP and Local Adam and remain roughly $2\times$ smaller than DDP thereafter, suggesting that Nesterov may steer optimization toward smoother regions of the loss landscape. State synchronization slightly accelerates the decrease in gradient norm over DiLoCo, after which the curves coincide.

Under DiLoCo, total output activation norms grow monotonically to more than $2\times$ the Adam-DDP values, whereas \texttt{DES-LOC} Nesterov substantially slows this growth, ending $\approx 32\%$ below DiLoCo (bottom row). This resembles the stabilization seen for \texttt{DES-LOC} versus FedAvg baselines without optimizer-state averaging (Fig.~5), where periodic averaging also curbed activation growth. This supports viewing finite synchronization as a regularizer that limits worker drift and optimizer-state noise, yielding better-controlled activations while offering Nesterov's benefits. \texttt{DES-LOC} Nesterov does incur additional communication costs relative to DiLoCo at fixed $K_x$, being $\approx 2\%$ slower than DiLoCo under these bandwidth conditions (Fig.~\ref{fig:nesterov_diloco_1b} bottom left) and $K_u, K_v$ settings while being $\approx 8\%$ faster than Local Adam. The extra cost can be made arbitrarily small by increasing the optimizer-state sync periods ($K_u, K_v$): in the limit $K_u, K_v \rightarrow \infty$ it recovers DiLoCo in both performance and communication, while any finite sync period partly inherits the robustness and fault-tolerance benefits of synchronizing optimizer states, modulated by the chosen $\beta$'s.

\begin{center}\fbox{\parbox{0.95\textwidth}{
\textbf{Takeaway:} At the 1B scale and long horizons, Nesterov-based local-update methods (DiLoCo, \texttt{DES-LOC} Nesterov) outperform Adam-DDP, consistent with prior scaling-law results. Relative to DiLoCo, \texttt{DES-LOC} Nesterov matches or improves perplexity while substantially reducing gradient and activation norms via periodic optimizer-state synchronization, yielding a tunable point on the communication--performance Pareto frontier.
}}\end{center}

%=============================================================================
\subsection{Very Low Bandwidth Experiments}
\label{sec:b6_low_bandwidth}
%=============================================================================

While perplexity is invariant to network bandwidth, wall-clock time is highly sensitive to it. To practically showcase this, we perform a benchmark with a 1B model to measure time under extremely low bandwidth conditions (10 Gbit/s). This setup simulates a scenario with affordable, consumer-grade interconnects rather than data-centers. Due to the extreme gradient synchronization delay inherent to DDP in this regime, the benchmark was limited to a $10{,}240$ step horizon to remain feasible.

\begin{figure}[t]
\centering
% Figure 20: 2-panel bandwidth benchmark (PDF page 28 top)
\caption{Training efficiency benchmark on a 1B model under restricted bandwidth (10 Gbit/s). \textbf{Left:} Perplexity versus wall-clock time. \texttt{DES-LOC} Nesterov effectively decouples training time from bandwidth, finishing in $\approx 9$ hours compared to the projected $\approx 3.5$ days for DDP (dashed blue). \textbf{Right:} Perplexity versus sequential steps. While step-wise convergence is comparable, the communication overhead of DDP creates a massive bottleneck in the time domain.}
\label{fig:low_bandwidth}
\end{figure}

As shown in Figure~\ref{fig:low_bandwidth}, \texttt{DES-LOC} Nesterov dramatically reduces training time by $\approx 9.42\times$ compared to DDP, completing the run in $8.99$ hours versus $84.73$ hours ($3.5$ days) for DDP, even with the constant overheads of our unoptimized implementation. Furthermore, \texttt{DES-LOC} Nesterov is more efficient than Local Adam, finishing $\approx 7\%$ faster (8.99h vs.\ 9.62h) while achieving significantly lower perplexity (8.91 vs.\ $10.19$). When compared to the ultra-lightweight DiLoCo baseline (8.68h), \texttt{DES-LOC} Nesterov incurs a time penalty of $\approx 3.6\%$ due to the additional optimizer state synchronization. However, this yields performance gains, improving final perplexity by $\approx 2.3\%$ (8.91 vs.\ 9.12) over DiLoCo.

\begin{center}\fbox{\parbox{0.95\textwidth}{
\textbf{Takeaway:} In extremely low bandwidth settings ($10$ Gbit/s), \texttt{DES-LOC} Nesterov eliminates the communication bottleneck, reducing training time by $9.42\times$ over DDP. It strikes a balance on the Pareto frontier: its wall-clock time is in-between those of Local Adam and DiLoCo while outperforming them both in perplexity.
}}\end{center}

%=============================================================================
\subsection{Muon Experimental Details (See Section~5.6)}
\label{sec:b7_muon_details}
%=============================================================================

This section provides the full experimental details for the Muon experiments presented in Section~5.6. Although a comprehensive theoretical treatment of preconditioned local updates under Newton-Schulz iterations is outside the scope of this work, the \texttt{DES-LOC} design is inherently compatible with such structures, as the momentum buffer follows standard EMA dynamics regardless of the subsequent preconditioning step.

\textbf{Experimental Details.} We utilize the standard PyTorch implementation of \texttt{Muon} with Nesterov momentum enabled and a weight decay of $0.1$. Following the recommendations of \citet{liu2025muon}, we apply the \texttt{match\_rms\_norm} adjustment to learning rates. We adopt the conventional split optimization strategy for \texttt{Muon}: AdamW handles embeddings and layer normalizations, while \texttt{Muon} optimizes all 2D matrices~\citep{jordan2024muon}. The momentum parameter for \texttt{Muon} is set to $\beta=0.9$, while the Adam component retains the $\beta_1=0.9, \beta_2=0.999$ settings used elsewhere. Gradient clipping thresholds are scaled by model size: $1.0$ for 16M, $0.5$ for 125M, and $0.25$ for 360M. For the \texttt{Local Muon} baseline, all optimizer states (Muon momentum; Adam first/second momenta) synchronize every $32$ steps. In contrast, \texttt{DES-LOC} delays state synchronization: the first momentum (for both optimizers) synchronizes every $96$ steps ($3\times$ reduction), and Adam's second momentum synchronizes every $192$ steps ($6\times$ reduction).

%=============================================================================
\subsection{Experiments on the Flux Vision Model}
\label{sec:b8_flux}
%=============================================================================

To demonstrate the universality of \texttt{DES-LOC} across different modalities and architectures beyond standard decoder-only LLMs, we evaluate its performance on \texttt{Flux}~\citep{labs2025flux}, a Rectified Flow Transformer designed for text-to-image generation. This architecture differs significantly from the causal language models evaluated in previous sections, serving as a robust test for the generalizability of our decoupled synchronization approach.

\textbf{Experimental Setup.} We utilize the $280$M parameter variant of \texttt{Flux} provided by \texttt{torchtitan}, training with a global batch size of $256$. The inner optimizer is AdamW with $\beta_1=0.9, \beta_2=0.999$. We compare three settings: (1)~DDP, (2)~Local Adam with a synchronization period of $K=32$, and (3)~\texttt{DES-LOC} with a parameter sync period of $K_x=32$. For \texttt{DES-LOC}, we decouple the momentum synchronization significantly, setting $K_u=3K_x$ and $K_v=6K_x$ (192 steps).

\begin{figure}[t]
\centering
% Figure 21: Flux training loss (PDF page 29 middle)
\caption{Training loss comparison on the $280$M parameter \texttt{Flux} model (Rectified Flow Transformer). \texttt{DES-LOC} ($K_x=32$, $K_u=192$, $K_v=192$) effectively matches the convergence trajectory of both the fully synchronous DDP baseline and Local Adam ($K=32$).}
\label{fig:flux}
\end{figure}

Our results, visualized in Fig.~\ref{fig:flux}, indicate that \texttt{DES-LOC} generally matches the performance of Local Adam and approaches the DDP upper bound.

\begin{center}\fbox{\parbox{0.95\textwidth}{
\textbf{Takeaway:} The efficacy of \texttt{DES-LOC} extends beyond LLMs to Rectified Flow Transformers (\texttt{Flux}). The method generally matches the performance of DDP and Local Adam while reducing communication by $2\times$ over Local Adam, demonstrating the universality of the approach. We leave the scaling of this result to larger vision models for future work.
}}\end{center}

%=============================================================================
\subsection{Throughput at 7B Scale}
\label{sec:b9_throughput}
%=============================================================================

To assess the practical scalability of our method on state-of-the-art hardware and at large model scales, we measure the training throughput of a 7B parameter model distributed across $8$ independent NVIDIA B200 GPUs.

\textbf{Throughput Analysis.} As illustrated in Fig.~\ref{fig:throughput_7b}, during the local update phases, each GPU operates at the peak efficiency of a fully isolated local run, achieving identical tokens-per-second throughput as a single B200 with zero synchronization overhead. Distinct drops in throughput are observed only at the sparse synchronization boundaries ($K_x=32$, $K_u=96$, $K_v=192$), where the system pauses to aggregate model parameters and optimizer states.

\begin{figure}[t]
\centering
% Figure 22: Throughput plot (PDF page 30 top)
\caption{Instantaneous throughput (tokens/sec) for a 7B model on 8x B200s. \texttt{DES-LOC} maintains peak ``local-only'' speed for the vast majority of steps, with throughput dips occurring only at synchronization intervals (32, 96, 192). In contrast, standard DDP would incur a synchronization penalty at \textit{every} step, permanently depressing the throughput curve.}
\label{fig:throughput_7b}
\end{figure}

Crucially, standard DDP incurs this communication penalty at \textit{every single training step}, significantly lowering the average tokens/sec. Even with our current unoptimized ``stop-the-world'' implementation---which explicitly pauses computation to communicate and does not leverage computation-communication overlap---\texttt{DES-LOC} significantly increases aggregate throughput by amortizing these costs over long local training windows.

\begin{center}\fbox{\parbox{0.95\textwidth}{
\textbf{Takeaway:} On high-performance B200 hardware, \texttt{DES-LOC} enables near-linear scaling by keeping workers in a high-throughput local regime for the majority of training. By restricting communication overhead to sparse intervals, it delivers significant wall-clock speedups over DDP, even without low-level implementation optimizations like communication overlap.
}}\end{center}

%=============================================================================
\section{Further Algorithmic Details of DES-LOC}
\label{sec:appendix_c}
%=============================================================================

\subsection{Extension to FedOpt}
\label{sec:fedopt}

Although~\citet{cheng2025convergence} show provable convergence for adaptive inner optimizers in a
federated optimization framework, their result rests on the assumption that after a period of local
work, the new global model is created by averaging the local client models. In relation to the larger
\texttt{FedOpt} literature~\citep{reddi2021adaptive}, the scheme chosen by~\citet{cheng2025convergence} resembles
that of \texttt{FedAvg}, or where the server optimizer is \texttt{SGD} with the outer learning rate set to one~\citep{reddi2021adaptive}.
Naturally, the question of whether alternate server optimizers than have been used in prior works
can also be implemented for \texttt{Local Adam}, and thus \texttt{DES-LOC}, arises.

We argue that indeed \texttt{DES-LOC}'s principles can be effectively applied to any \texttt{FedOpt} method and
not just \texttt{FedAvg}. While using an alternate server optimizer does not have proven convergence
guarantees as yet, we show in Algorithm~\ref{alg:desloc} that the choice of the \texttt{ServerOpt} is not constrained
from a practical point-of-view. However, the improvements that \texttt{DES-LOC} provide are related to
the local optimization procedure, which is orthogonal to the outer optimizer choice.
Choosing the correct, and most effective, outer optimizer is an open research area~\citep{khaled2025tuning},
and we leave the investigations of the interactions between \texttt{DES-LOC} and outer optimizers to future work.

\subsection{Deterministic Optimizer-specific Variants of Algorithm 1}
\label{sec:algorithm_variants}

\begin{algorithm}[h]
\caption{DES-LOC-Adam}
\label{alg:desloc_adam}
\begin{algorithmic}[1]
\REQUIRE Model tensors, Hyper-parameters
\STATE $x_0, u_{-1}, v_{-1} \in \mathbb{R}^d$ \COMMENT{initial parameter vector, seeds for first and second moments}
\STATE $\{\eta_t\}_{t=0}^{T-1} \subset \mathbb{R}_{>0}$ \COMMENT{step-size schedule}
\STATE $\beta_1, \beta_2 \in [0,1)$ \COMMENT{Adam decay factors}
\STATE $\rho, \lambda \in \mathbb{R}_{>0}$ \COMMENT{gradient clipping term, $\ell_2$ stability term}
\STATE $T, M \in \mathbb{N}_+$ \COMMENT{total iterations, number of workers}
\STATE $K_x, K_u, K_v \in \mathbb{N}_+$ \COMMENT{sync periods for parameters, first and second moments}
\ENSURE $x_T, u_{T-1}, v_{T-1}$
\STATE for each worker $m$: $x^m_0 = x_0$, $u^m_{-1} = v^m_{-1} = 0$ \COMMENT{local init ($t = -1$ seeds training loop)}
\FOR{$t = 0, \ldots, T-1$} \COMMENT{training loop}
    \FOR{all workers $m = 0, \ldots, M-1$ \textbf{in parallel}}
        \STATE $g^m_t \gets \nabla F(x^m_t; \xi^m_t)$ \COMMENT{stochastic gradient}
        \STATE $\hat{g}^m_t \gets \text{clip}(g^m_t, \rho)$ \COMMENT{clip to radius $\rho$}
        \IF{$t \bmod K_u = 0$} \COMMENT{sync $u$}
            \STATE $u^m_t \gets \beta_1 \mathbb{E}_m[u^m_{t-1}] + (1 - \beta_1)\hat{g}^m_t$
        \ELSE
            \STATE $u^m_t \gets \beta_1 u^m_{t-1} + (1 - \beta_1)\hat{g}^m_t$
        \ENDIF
        \IF{$t \bmod K_v = 0$} \COMMENT{sync $v$}
            \STATE $v^m_t \gets \beta_2 \mathbb{E}_m[v^m_{t-1}] + (1 - \beta_2)(\hat{g}^m_t \odot \hat{g}^m_t)$
        \ELSE
            \STATE $v^m_t \gets \beta_2 v^m_{t-1} + (1 - \beta_2)(\hat{g}^m_t \odot \hat{g}^m_t)$
        \ENDIF
        \STATE $d^m_t \gets \frac{\eta_t}{\sqrt{v^m_t + \lambda^2}} \odot u^m_t$ \COMMENT{bias-corrected step}
        \IF{$t \bmod K_x = 0$} \COMMENT{sync $x$}
            \STATE $x^m_{t+1} \gets \mathbb{E}_m[x^m_t] - d^m_t$
        \ELSE
            \STATE $x^m_{t+1} \gets x^m_t - d^m_t$
        \ENDIF
    \ENDFOR
\ENDFOR
\end{algorithmic}
\end{algorithm}

\begin{algorithm}[h]
\caption{DES-LOC-ADOPT}
\label{alg:desloc_adopt}
\begin{algorithmic}[1]
\REQUIRE Model tensors, Hyper-parameters
\STATE $x_0, m_{-1}, v_{-1} \in \mathbb{R}^d$ \COMMENT{initial parameter vector and momenta}
\STATE $\{\eta_t\}_{t=0}^{T-1} \subset \mathbb{R}_{>0}$ \COMMENT{learning rate schedule}
\STATE $\beta_1, \beta_2 \in [0,1)$ \COMMENT{decay factors}
\STATE $\rho, \epsilon \in \mathbb{R}_{>0}$ \COMMENT{gradient clipping term, small stability constant}
\STATE $T, M \in \mathbb{N}_+$ \COMMENT{total iterations, number of workers}
\STATE $K_x, K_m, K_v \in \mathbb{N}_+$ \COMMENT{sync periods for parameters, first and second moments}
\ENSURE $x_T$, $m_{T-1}$, $v_{T-1}$
\STATE for each worker $m$: $x^m_0 = x_0$, $m^m_{-1} = v^m_{-1} = 0$ \COMMENT{local initialization}
\FOR{$t = 0, \ldots, T-1$}
    \FOR{all workers $m = 0, \ldots, M-1$ \textbf{in parallel}}
        \STATE $g^m_t \gets \nabla F(x^m_t; \xi^m_t)$ \COMMENT{stochastic gradient}
        \STATE $\hat{g}^m_t \gets \text{clip}(g^m_t, \rho)$ \COMMENT{gradient clipping}
        \IF{$t \bmod K_v = 0$}
            \STATE $v^m_t \gets \beta_2 \mathbb{E}_m[v^m_{t-1}] + (1 - \beta_2)(\hat{g}^m_t \odot \hat{g}^m_t)$
        \ELSE
            \STATE $v^m_t \gets \beta_2 v^m_{t-1} + (1 - \beta_2)(\hat{g}^m_t \odot \hat{g}^m_t)$
        \ENDIF
        \IF{$t \bmod K_m = 0$}
            \STATE $m^m_t \gets \beta_1 \mathbb{E}_m[m^m_{t-1}] + (1 - \beta_1) \frac{\hat{g}^m_t}{\max\{\sqrt{v^m_{t-1}}, \epsilon\}}$
        \ELSE
            \STATE $m^m_t \gets \beta_1 m^m_{t-1} + (1 - \beta_1) \frac{\hat{g}^m_t}{\max\{\sqrt{v^m_{t-1}}, \epsilon\}}$ \COMMENT{ADOPT update}
        \ENDIF
        \STATE $d^m_t \gets \eta_t m^m_t$
        \IF{$t \bmod K_x = 0$}
            \STATE $x^m_{t+1} \gets \mathbb{E}_m[x^m_t] - d^m_t$
        \ELSE
            \STATE $x^m_{t+1} \gets x^m_t - d^m_t$
        \ENDIF
    \ENDFOR
\ENDFOR
\end{algorithmic}
\end{algorithm}

%=============================================================================
% Appendix D: Convergence Analysis of DES-LOC-SGDM
% Claude-12 (M276-M300): PDF pages 32-35 (SGDM proof + D.1 Adam extension start)
%=============================================================================
\section{Convergence Analysis of DES-LOC-SGDM (in expectation bounds)}
\label{sec:sgdm_proof}

Here we provide a non-convex convergence analysis of the proposed \texttt{DES-LOC} approach applied to the SGDM optimizer which has a single state ($N=1$, momentum). The complete description of the algorithm can be found in Algorithm~\ref{alg:desloc_sgdm}.

\begin{algorithm}[ht]
\caption{\texttt{DES-LOC-SGDM}}
\label{alg:desloc_sgdm}
\begin{algorithmic}[1]
\REQUIRE \textbf{Model tensors}
\STATE $x_0 \in \mathbb{R}^d$ --- initial parameter vector
\STATE $u_{-1} \in \mathbb{R}^d$ --- seed for the momentum, initialised to $\mathbf{0}$
\REQUIRE \textbf{Hyper-parameters}
\STATE $\{\eta_t\}_{t=0}^{T-1} \subset \mathbb{R}_{>0}$ --- step-size schedule
\STATE $\beta \in [0,1)$ --- Momentum decay factor
\STATE $T \in \mathbb{N}_+$ --- total optimisation iterations
\STATE $M \in \mathbb{N}_+$ --- number of workers
\STATE $p_x = \frac{1}{K_x},\; p_u = \frac{1}{K_u} \in [0,1]$ --- synchronization probabilities for parameters and momentum
\ENSURE $x_T$, $u_{T-1}^m$
\STATE \textbf{for each worker} $m$: $x_0^m = x_0,\; u_{-1}^m = 0$ \hfill local init ($t=-1$ seeds)
\FOR{$t = 0, \ldots, T-1$} \hfill training loop
    \FOR{all workers $m = 0, \ldots, M-1$ \textbf{in parallel}}
        \STATE $g_t^m \leftarrow \nabla F_m(x_t^m; \xi_t^m)$ \hfill stochastic gradient
        \STATE $u_t^m \leftarrow \begin{cases} \mathbb{E}_m[\beta u_{t-1}^m + (1-\beta) g_t^m], & \text{with probability } p_u \\ \beta u_{t-1}^m + (1-\beta) g_t^m, & \text{with probability } 1-p_u \end{cases}$ \hfill sync $u$
        \STATE $x_{t+1}^m \leftarrow \begin{cases} \mathbb{E}_m[x_t^m - \eta_t u_t^m], & \text{with probability } p_x \\ x_t^m - \eta_t u_t^m, & \text{with probability } 1-p_x \end{cases}$ \hfill sync $x$
    \ENDFOR
\ENDFOR
\end{algorithmic}
\end{algorithm}

In order to facilitate the technical presentation, we model synchronization frequencies by assigning probabilities to each averaging event. For example, the parameters $x_t^m$ are synchronized with the probability $p_x = \frac{1}{K_x}$, which is statistically equivalent to performing the averaging in every $\frac{1}{p_x} = K_x$ iteration. Similarly, momentum $u_t^m$ synchronization happens with probability $p_u = \frac{1}{K_u}$, which can differ from $p_x$.

\paragraph{Step 1 (virtual iterates).} For each step $t \geq 0$, denote the average parameters, momentum and gradient as follows:
\[
x_t \stackrel{\text{def}}{=} \mathbb{E}_m[x_t^m], \quad
u_t \stackrel{\text{def}}{=} \mathbb{E}_m[u_t^m], \quad
g_t \stackrel{\text{def}}{=} \mathbb{E}_m[g_t^m].
\]
Then these averaged variables follow the ``standard'' centralized SGDM dynamics:
\begin{align*}
u_t &= \beta u_{t-1} + (1-\beta) g_t \\
x_{t+1} &= x_t - \eta u_t.
\end{align*}
Letting $x_{-1} = x_0$, define the global virtual iterations as follows
\[
z_t \stackrel{\text{def}}{=} \frac{1}{1-\beta} x_t - \frac{\beta}{1-\beta} x_{t-1}, \quad t \geq 0.
\]
The key property of this virtual iterates we are going to exploit in the next steps is that they follow averaged gradients, namely for any $t \geq 0$ we have
\begin{align*}
z_{t+1} - z_t &= \frac{1}{1-\beta}(x_{t+1} - x_t) - \frac{\beta}{1-\beta}(x_t - x_{t-1}) \\
&= -\frac{\eta}{1-\beta} u_t + \frac{\eta\beta}{1-\beta} u_{t-1}
 = -\frac{\eta}{1-\beta}(u_t - \beta u_{t-1}) = -\eta g_t.
\end{align*}

\paragraph{Step 2 (smoothness over virtual iterates).} Then we apply smoothness of the global loss function $f$ over these global virtual iterates.
\begin{align*}
f(z_{t+1}) &\leq f(z_t) + \langle \nabla f(z_t), z_{t+1} - z_t \rangle + \frac{L}{2}\|z_{t+1} - z_t\|^2 \\
&= f(z_t) + \langle \underbrace{\nabla f(x_t), z_{t+1} - z_t}_{I} \rangle
  + \langle \underbrace{\nabla f(z_t) - \nabla f(x_t), z_{t+1} - z_t}_{II} \rangle
  + \underbrace{\frac{L}{2}\|z_{t+1} - z_t\|^2}_{III}.
\end{align*}
In the next step, we separately bound each term appearing in the above bound.

\paragraph{Step 3a (one step progress).} Bounding term I.
\begin{align*}
&\mathbb{E}\langle \nabla f(x_t), z_{t+1} - z_t \rangle \\
&= -\eta \mathbb{E}\left\langle \nabla f(x_t), \frac{1}{M}\sum_{m=1}^{M} g_t^m \right\rangle
 = -\eta \mathbb{E}\left\langle \nabla f(x_t), \frac{1}{M}\sum_{m=1}^{M} \nabla f_m(x_t^m) \right\rangle \\
&= -\frac{\eta}{2}\mathbb{E}\|\nabla f(x_t)\|^2
  -\frac{\eta}{2}\mathbb{E}\left\|\frac{1}{M}\sum_{m=1}^{M}\nabla f_m(x_t^m)\right\|^2
  +\frac{\eta}{2}\mathbb{E}\left\|\nabla f(x_t) - \frac{1}{M}\sum_{m=1}^{M}\nabla f_m(x_t^m)\right\|^2 \\
&= -\frac{\eta}{2}\mathbb{E}\|\nabla f(x_t)\|^2
  -\frac{\eta}{2}\mathbb{E}\left\|\frac{1}{M}\sum_{m=1}^{M}\nabla f_m(x_t^m)\right\|^2
  +\frac{\eta}{2M}\sum_{m=1}^{M}\mathbb{E}\|\nabla f_m(x_t) - \nabla f_m(x_t^m)\|^2 \\
&\leq -\frac{\eta}{2}\mathbb{E}\|\nabla f(x_t)\|^2
  -\frac{\eta}{2}\mathbb{E}\left\|\frac{1}{M}\sum_{m=1}^{M}\nabla f_m(x_t^m)\right\|^2
  +\frac{\eta L^2}{2M}\sum_{m=1}^{M}\underbrace{\mathbb{E}\|x_t - x_t^m\|^2}_{\text{Lemma~\ref{lem:bound}}}.
\end{align*}

\paragraph{Step 3b (one step progress).} Bounding term II.
\begin{align*}
&\mathbb{E}\langle \nabla f(z_t) - \nabla f(x_t), z_{t+1} - z_t \rangle
= -\eta \mathbb{E}\left\langle \nabla f(z_t) - \nabla f(x_t), \frac{1}{M}\sum_{m=1}^{M}\nabla f_m(x_t^m) \right\rangle \\
&\leq \frac{\eta\rho}{2}\mathbb{E}\|\nabla f(z_t) - \nabla f(x_t)\|^2
  + \frac{\eta}{2\rho}\mathbb{E}\left\|\frac{1}{M}\sum_{m=1}^{M}\nabla f_m(x_t^m)\right\|^2 \\
&\leq \frac{\eta\rho L^2}{2}\underbrace{\mathbb{E}\|z_t - x_t\|^2}_{\text{Lemma~\ref{lem:virt_drift}}}
  + \frac{\eta}{2\rho}\mathbb{E}\left\|\frac{1}{M}\sum_{m=1}^{M}\nabla f_m(x_t^m)\right\|^2.
\end{align*}

\paragraph{Step 3c (one step progress).} Bounding term III.
\begin{align*}
\frac{L}{2}\mathbb{E}\|z_{t+1} - z_t\|^2
&= \frac{\eta^2 L}{2}\mathbb{E}\left\|\frac{1}{M}\sum_{m=1}^{M} g_t^m\right\|^2 \\
&= \frac{\eta^2 L}{2}\mathbb{E}\left\|\frac{1}{M}\sum_{m=1}^{M} g_t^m - \nabla f_m(x_t^m)\right\|^2
  + \frac{\eta^2 L}{2}\mathbb{E}\left\|\frac{1}{M}\sum_{m=1}^{M}\nabla f_m(x_t^m)\right\|^2 \\
&= \frac{\eta^2 L}{2M^2}\sum_{m=1}^{M}\mathbb{E}\|g_t^m - \nabla f_m(x_t^m)\|^2
  + \frac{\eta^2 L}{2}\mathbb{E}\left\|\frac{1}{M}\sum_{m=1}^{M}\nabla f_m(x_t^m)\right\|^2 \\
&\leq \frac{\eta^2 L}{2M}\sigma^2
  + \frac{\eta^2 L}{2}\mathbb{E}\left\|\frac{1}{M}\sum_{m=1}^{M}\nabla f_m(x_t^m)\right\|^2.
\end{align*}

\paragraph{Step 3abc (one step progress).} Combining previous bounds.
\begin{align*}
\mathbb{E}f(z_{t+1}) - \mathbb{E}f(z_t)
&\leq \mathbb{E}\underbrace{\langle \nabla f(x_t), z_{t+1} - z_t \rangle}_{I}
  + \mathbb{E}\underbrace{\langle \nabla f(z_t) - \nabla f(x_t), z_{t+1} - z_t \rangle}_{II}
  + \mathbb{E}\underbrace{\frac{L}{2}\|z_{t+1} - z_t\|^2}_{III} \\
&\leq -\frac{\eta}{2}\mathbb{E}\|\nabla f(x_t)\|^2
  -\frac{\eta}{2}\mathbb{E}\left\|\frac{1}{M}\sum_{m=1}^{M}\nabla f_m(x_t^m)\right\|^2
  +\frac{\eta L^2}{2M}\sum_{m=1}^{M}\underbrace{\mathbb{E}\|x_t - x_t^m\|^2}_{\text{Lemma~\ref{lem:bound}}} \\
&\quad +\frac{\eta\rho L^2}{2}\underbrace{\mathbb{E}\|z_t - x_t\|^2}_{\text{Lemma~\ref{lem:virt_drift}}}
  +\frac{\eta}{2\rho}\mathbb{E}\left\|\frac{1}{M}\sum_{m=1}^{M}\nabla f_m(x_t^m)\right\|^2 \\
&\quad +\frac{\eta^2 L}{2M}\sigma^2
  + \frac{\eta^2 L}{2}\mathbb{E}\left\|\frac{1}{M}\sum_{m=1}^{M}\nabla f_m(x_t^m)\right\|^2.
\end{align*}

\paragraph{Step 4 (final).} Now we average over the iterates and apply the bounds derived in Lemmas~\ref{lem:virt_drift},~\ref{lem:bound}.
\begin{align*}
\frac{\mathbb{E}[f(z_T) - f(z_0)]}{T}
&= \frac{1}{T}\sum_{t=0}^{T-1}\mathbb{E}[f(z_{t+1}) - f(z_t)] \\
&\leq -\frac{\eta}{2T}\sum_{t=0}^{T-1}\mathbb{E}\|\nabla f(x_t)\|^2
  - \frac{\eta}{2}\left(1 - \frac{1}{\rho} - \eta L\right)
    \frac{1}{T}\sum_{t=0}^{T-1}\mathbb{E}\left\|\frac{1}{M}\sum_{m=1}^{M}\nabla f_m(x_t^m)\right\|^2 \\
&\quad + \frac{\eta\rho L^2}{2}\underbrace{\frac{1}{T}\sum_{t=0}^{T-1}\mathbb{E}\|z_t - x_t\|^2}_{\text{Lemma~\ref{lem:virt_drift}}}
  + \frac{\eta L^2}{2M}\underbrace{\frac{1}{TM}\sum_{t=0}^{T-1}\sum_{m=1}^{M}\mathbb{E}\|x_t - x_t^m\|^2}_{\text{Lemma~\ref{lem:bound}}}
  + \frac{\eta^2 L}{2M}\sigma^2 \\
&\leq -\frac{\eta}{2T}\sum_{t=0}^{T-1}\mathbb{E}\|\nabla f(x_t)\|^2
  - \frac{\eta}{2}\left(1 - \frac{1}{\rho} - \eta L\right)
    \frac{1}{T}\sum_{t=0}^{T-1}\mathbb{E}\left\|\frac{1}{M}\sum_{m=1}^{M}\nabla f_m(x_t^m)\right\|^2
  + \frac{\eta^2 L}{2M}\sigma^2 \\
&\quad + \frac{\eta\rho L^2}{2}\left(
    \frac{\eta^2\beta^2}{(1-\beta)^2 M}\sigma^2
    + \frac{\eta^2\beta^2}{(1-\beta)^2}\frac{1}{T}\sum_{\tau=0}^{T-1}\mathbb{E}\left\|\frac{1}{M}\sum_{m=1}^{M}\nabla f_m(x_\tau^m)\right\|^2
  \right) \\
&\quad + \frac{\eta L^2}{2}\left(
    12\eta^2(B^2-1)\psi \cdot \frac{1}{T}\sum_{t=0}^{T-1}\mathbb{E}\|\nabla f(x_t)\|^2
    + 4\eta^2\psi(\sigma^2 + 3G^2)
  \right) \\
&\leq -\frac{\eta}{2}\left(1 - 12\eta^2 L^2(B^2-1)\psi\right)
    \frac{1}{T}\sum_{t=0}^{T-1}\mathbb{E}\|\nabla f(x_t)\|^2 \\
&\quad - \frac{\eta}{2}\left(1 - \frac{1}{\rho} - \eta L - \frac{\eta^2\beta^2\rho L^2}{(1-\beta)^2}\right)
    \frac{1}{T}\sum_{t=0}^{T-1}\mathbb{E}\left\|\frac{1}{M}\sum_{m=1}^{M}\nabla f_m(x_t^m)\right\|^2 \\
&\quad + \frac{\eta^2 L}{2M}\sigma^2
  + \frac{\eta^3\rho L^2\beta^2}{2(1-\beta)^2 M}\sigma^2
  + 2\eta^3 L^2\psi(\sigma^2 + 3G^2).
\end{align*}

Next, we choose $\rho = 2$ and step size $\eta$ such that
\begin{align*}
12\eta^2 L^2(B^2-1)\psi &\leq \frac{1}{2} & &\Longleftrightarrow \quad \text{to bound the first term} \\
\eta L + \frac{2\eta^2\beta^2 L^2}{(1-\beta)^2} &\leq \frac{1}{2} & &\Longleftrightarrow \quad \text{to bound the second term} \\
12\eta^2 L^2\psi &\leq \frac{1}{2} & &\Longleftrightarrow \quad \text{from Lemma~\ref{lem:bound}}
\end{align*}
Note that
\[
\eta_0 \stackrel{\text{def}}{=} \frac{1}{4L}\min\!\left(1-\beta,\; \frac{1}{6\sqrt{\psi\max(1,B^2-1)}}\right)
\]
satisfies all three bounds. Then, with any $\eta \leq \eta_0$ we get
\begin{align*}
\frac{\mathbb{E}[f(z_T) - f(z_0)]}{T}
&\leq -\frac{\eta}{4T}\sum_{t=0}^{T-1}\mathbb{E}\|\nabla f(x_t)\|^2 \\
&\quad + \frac{\eta^2 L}{2M}\sigma^2
  + \frac{\eta^3\rho L^2\beta^2}{2(1-\beta)^2 M}\sigma^2
  + 2\eta^3 L^2\psi(\sigma^2 + 3G^2).
\end{align*}
Noticing that $z_0 = x_0$ and $f^* \leq f(z_T)$, we have
\[
\frac{1}{T}\sum_{t=0}^{T-1}\mathbb{E}\|\nabla f(x_t)\|^2
\leq \frac{4(f(x_0) - f^*)}{\eta T}
  + \frac{2\eta L}{M}\sigma^2
  + \frac{4\eta^2 L^2\beta^2}{(1-\beta)^2 M}\sigma^2
  + 8\eta^2 L^2\psi(\sigma^2 + 3G^2).
\]
Furthermore, choosing $\eta = \min(\eta_0, \frac{1}{\sqrt{T}})$, we get the following rate:
\begin{align*}
&\frac{1}{T}\sum_{t=0}^{T-1}\mathbb{E}\|\nabla f(x_t)\|^2 \\
&\leq \max\!\left(1, \frac{1}{\eta_0\sqrt{T}}\right)\frac{4(f(x_0)-f^*)}{\sqrt{T}}
  + \frac{2L\sigma^2}{M\sqrt{T}}
  + \frac{4L^2\beta^2\sigma^2}{(1-\beta)^2 MT}
  + \frac{8L^2\psi(\sigma^2+3G^2)}{T} \\
&\leq \frac{4(f(x_0)-f^*)}{\sqrt{T}}
  + \frac{2L\sigma^2}{M\sqrt{T}}
  + \frac{4(f(x_0)-f^*)}{\eta_0 T}
  + \frac{4L^2\beta^2\sigma^2}{(1-\beta)^2 MT}
  + \frac{8L^2\psi(\sigma^2+3G^2)}{T} \\
&= \frac{4}{\sqrt{T}}\left(f(x_0) - f^* + \frac{L\sigma^2}{2M}\right)
  + \mathcal{O}\!\left(\frac{1+\psi}{T}\right).
\end{align*}

%-----------------------------------------------------------------------------
% D.1 Extension to Adam optimizer (PDF page 35 bottom)
%-----------------------------------------------------------------------------
\subsection{Extension to Adam optimizer}
\label{sec:adam_extension}

Here we discuss extension of the previous analysis for the Adam optimizer including the second-order momentum in the analysis. The addition is similar to the first-order momentum while the synchronization probability $p_v$ can differ from other probabilities $p_x$ and $p_u$. The complete description of the algorithm can be found in Algorithm~\ref{alg:desloc_adam_prob}. Instead of bounded heterogeneity Assumption~3, in this analysis we use stronger condition mentioned below:

\begin{assumption}[Bounded gradient]
\label{ass:bounded_grad}
For any iterate $t \geq 0$ and worker $m$, the local stochastic gradient is bounded, namely $\|g_t^m\|_2 \leq G$.
\end{assumption}

This condition facilitates the analysis by providing uniform upper bounds for gradients/momentum variables and is commonly used in the analysis of adaptive optimization.

\paragraph{Step 1 (preconditioning and virtual iterates).} Let $\Gamma_t^m \stackrel{\text{def}}{=} \mathbf{diag}^{-1/2}(\tilde{v}_t^m + \lambda^2)$ be the preconditioning matrix and for each step $t \geq 0$, denote the averaged variables
\[
x_t \stackrel{\text{def}}{=} \mathbb{E}_m[x_t^m], \quad
u_t \stackrel{\text{def}}{=} \mathbb{E}_m[u_t^m], \quad
v_t \stackrel{\text{def}}{=} \mathbb{E}_m[v_t^m], \quad
\tilde{v}_t \stackrel{\text{def}}{=} \mathbb{E}_m[\tilde{v}_t^m], \quad
g_t \stackrel{\text{def}}{=} \mathbb{E}_m[g_t^m].
\]
Then
\begin{align*}
u_t &= \beta_1 u_{t-1} + (1-\beta_1) g_t \\
x_{t+1} &= x_t - d_t = x_t - \eta\,\mathbb{E}_m[\Gamma_t^m u_t^m].
\end{align*}

% Claude-13 (M301-M325): PDF pages 36-43 (D.1 Adam proof cont. + Key Lemmas 2,3,4)

%-----------------------------------------------------------------------------
% Algorithm 5: DES-LOC-Adam (probabilistic synchronization) — PDF page 36 top
%-----------------------------------------------------------------------------
\begin{algorithm}[h]
\caption{\texttt{DES-LOC-Adam} (with probabilistic synchronization)}
\label{alg:desloc_adam_prob}
\begin{algorithmic}[1]
\REQUIRE \textbf{Model tensors}
\STATE $x_0 \in \mathbb{R}^d$ --- initial parameter vector
\STATE $u_{-1}, v_{-1} \in \mathbb{R}^d$ --- seeds for first and second moments, initialised to $\mathbf{0}$
\REQUIRE \textbf{Hyper-parameters}
\STATE $\{\eta_t\}_{t=0}^{T-1} \subset \mathbb{R}_{>0}$ --- step-size schedule
\STATE $\beta_1, \beta_2 \in [0,1)$ --- Adam decay factors
\STATE $\lambda \in \mathbb{R}_{\geq 0}$ --- $\ell_2$ stability term
\STATE $T \in \mathbb{N}_+$ --- total optimisation iterations
\STATE $M \in \mathbb{N}_+$ --- number of workers
\STATE $p_x = \frac{1}{K_x},\, p_u = \frac{1}{K_u},\, p_v = \frac{1}{K_v} \in [0,1]$ --- synchronization probabilities for parameters and momentums
\ENSURE $x_T$, $u_{T-1}$, $v_{T-1}$
\STATE \textbf{for each worker} $m$: $x_0^m = x_0$, $u_{-1}^m = v_{-1}^m = 0$ \COMMENT{local init ($t=-1$ seeds training loop)}
\FOR{$t = 0, \ldots, T-1$}
\FOR{all workers $m = 0, \ldots, M-1$ \textbf{in parallel}}
\STATE $g_t^m \leftarrow \nabla F_m(x_t^m; \xi_t^m)$ \COMMENT{stochastic gradient}
\STATE $u_t^m \leftarrow \begin{cases} \mathbb{E}_m[\beta_1 u_{t-1}^m + (1-\beta_1) g_t^m], & \text{with probability } p_u \\ \beta_1 u_{t-1}^m + (1-\beta_1) g_t^m, & \text{with probability } 1 - p_u \end{cases}$ \COMMENT{sync $u$}
\STATE $v_t^m \leftarrow \begin{cases} \mathbb{E}_m[\beta_2 v_{t-1}^m + (1-\beta_2)(g_t^m \odot g_t^m)], & \text{with probability } p_v \\ \beta_2 v_{t-1}^m + (1-\beta_2)(g_t^m \odot g_t^m), & \text{with probability } 1 - p_v \end{cases}$ \COMMENT{sync $v$}
\STATE $\tilde{v}_t^m \leftarrow \max(v_t^m, \tilde{v}_{t-1}^m)$ \COMMENT{AMSGrad Normalization, $\tilde{v}_{-1} = v_{-1}$}
\STATE $d_t^m \leftarrow \frac{\eta_t}{\sqrt{\tilde{v}_t^m + \lambda^2}} \odot u_t^m$ \COMMENT{bias-corrected update}
\STATE $x_{t+1}^m \leftarrow \begin{cases} \mathbb{E}_m[x_t^m - d_t^m], & \text{with probability } p_x \\ x_t^m - d_t^m, & \text{with probability } 1 - p_x \end{cases}$ \COMMENT{sync $x$}
\ENDFOR
\ENDFOR
\end{algorithmic}
\end{algorithm}

%-----------------------------------------------------------------------------
% Virtual iterates z_t for Adam (PDF page 36 bottom)
%-----------------------------------------------------------------------------
Consider the same averaged iterates $x_t$ and virtual iterates $z_t$ as before:
\[
z_t = \frac{1}{1-\beta_1} x_t - \frac{\beta_1}{1-\beta_1} x_{t-1}.
\]
In particular, $z_0 = x_0$. Then,
\begin{align*}
z_{t+1} - z_t
&= \frac{1}{1-\beta_1}(x_{t+1}-x_t) - \frac{\beta_1}{1-\beta_1}(x_t - x_{t-1}) \\
&= -\frac{\eta}{1-\beta_1}\mathbb{E}_m[\Gamma_t^m u_t^m]
   + \frac{\eta\beta_1}{1-\beta_1}\mathbb{E}_m[\Gamma_{t-1}^m u_{t-1}^m] \\
&= -\frac{\eta}{1-\beta_1}\mathbb{E}_m[\Gamma_t^m u_t^m]
   + \frac{\eta\beta_1}{1-\beta_1}\mathbb{E}_m[\Gamma_{t-1}^m u_{t-1}^m]
   \pm \frac{\eta\beta_1}{1-\beta_1}\mathbb{E}_m[\Gamma_t^m u_{t-1}^m] \\
&= -\frac{\eta}{1-\beta_1}\mathbb{E}_m[\Gamma_t^m(u_t^m - \beta_1 u_{t-1}^m)]
   + \frac{\eta\beta_1}{1-\beta_1}\mathbb{E}_m[(\Gamma_{t-1}^m - \Gamma_t^m) u_{t-1}^m] \\
&= -\eta\,\mathbb{E}_m[\Gamma_t^m \widetilde{g}_t^m]
   + \frac{\eta\beta_1}{1-\beta_1}\mathbb{E}_m[(\Gamma_{t-1}^m - \Gamma_t^m) u_{t-1}^m] \\
&= -\eta\,\mathbb{E}_m[\Gamma_t^m g_t]
   + \eta\,\mathbb{E}_m[\Gamma_t^m(g_t - \widetilde{g}_t^m)]
   + \frac{\eta\beta_1}{1-\beta_1}\mathbb{E}_m[(\Gamma_{t-1}^m - \Gamma_t^m) u_{t-1}^m] \\
&= -\eta\,\Gamma_t g_t
   + \eta \cdot \underbrace{\mathbb{E}_m[\Gamma_t^m(g_t - \widetilde{g}_t^m)]}_{\stackrel{\text{def}}{=}\, U_t}
   + \eta \cdot \underbrace{\frac{\beta_1}{1-\beta_1}\mathbb{E}_m[(\Gamma_{t-1}^m - \Gamma_t^m) u_{t-1}^m]}_{\stackrel{\text{def}}{=}\, V_t},
\end{align*}
where $\Gamma_t \stackrel{\text{def}}{=} \mathbb{E}_m[\Gamma_t^m]$ and $\widetilde{g}_t^m \stackrel{\text{def}}{=} \frac{u_t^m - \beta_1 u_{t-1}^m}{1-\beta_1}$ for which, $\mathbb{E}_m[\widetilde{g}_t^m] = \mathbb{E}_m[g_t^m] = g_t$.

%-----------------------------------------------------------------------------
% Step 2 (smoothness over virtual iterates for Adam) — PDF page 37 top
%-----------------------------------------------------------------------------
\paragraph{Step 2 (smoothness over virtual iterates).} Then we apply smoothness of the global loss function $f$ over these global virtual iterates.
\begin{align*}
f(z_{t+1}) - f(z_t)
&\leq \langle \nabla f(z_t), z_{t+1}-z_t \rangle + \frac{L}{2}\|z_{t+1}-z_t\|^2 \\
&= -\eta\langle \nabla f(z_t), \Gamma_t g_t \rangle
   + \eta\langle \nabla f(z_t), U_t \rangle
   + \eta\langle \nabla f(z_t), V_t \rangle
   + \frac{L}{2}\|z_{t+1}-z_t\|^2 \\
&= \underbrace{-\eta\langle \nabla f(x_t), \Gamma_t g_t \rangle}_{I}
   + \underbrace{\eta\langle \nabla f(z_t), U_t \rangle}_{II}
   + \underbrace{\eta\langle \nabla f(z_t), V_t \rangle}_{III} \\
&\quad + \underbrace{\frac{\eta^2 L}{2}\|\Gamma_t g_t - U_t - V_t\|^2}_{IV}
   + \underbrace{\eta\langle \nabla f(x_t) - \nabla f(z_t), \Gamma_t g_t \rangle}_{V}.
\end{align*}

In the next step, we separately bound each term appearing in the above bound. For clarity, we are also going to use $\|\nabla f(x_t)\| \leq G$ and $\|\nabla f(z_t)\| \leq G$. However, these conditions can be avoided through linking $\nabla f(z_t)$ term to $\nabla f(x_t)$, and $\nabla f(x_t)$ term to $\mathbb{E}_m \nabla f_m(x_t^m)$ with the bound for $\mathbb{E}[\|x_t - x_t^m\|^2]$.

%-----------------------------------------------------------------------------
% Step 3a (term I) — PDF page 37
%-----------------------------------------------------------------------------
\paragraph{Step 3a (one step progress). Bounding term I.}
\begin{align*}
I &= -\eta\langle \nabla f(x_t), \Gamma_t g_t \rangle \\
&= -\eta\,\mathbb{E}\!\left[\langle \nabla f(x_t), \Gamma_{t-1} g_t \rangle\right]
   + \eta\,\mathbb{E}\!\left[\langle \nabla f(x_t), (\Gamma_{t-1}-\Gamma_t) g_t \rangle\right] \\
&\leq -\eta\,\mathbb{E}\!\left[\left\langle \nabla f(x_t),\, \frac{1}{M}\sum_{m=1}^{M}\nabla f_m(x_t^m) \right\rangle_{\!\Gamma_{t-1}}\right]
   + \eta G^2 \mathbb{E}[\|\Gamma_{t-1}-\Gamma_t\|] \\
&\leq -\frac{\eta}{2}\mathbb{E}\!\left[\|\nabla f(x_t)\|_{\Gamma_{t-1}}^2\right]
   - \frac{\eta}{2}\mathbb{E}\!\left[\left\|\frac{1}{M}\sum_{m=1}^{M}\nabla f_m(x_t^m)\right\|_{\Gamma_{t-1}}^2\right] \\
&\quad + \frac{\eta}{2}\mathbb{E}\!\left[\left\|\nabla f(x_t) - \frac{1}{M}\sum_{m=1}^{M}\nabla f_m(x_t^m)\right\|_{\Gamma_{t-1}}^2\right]
   + \eta G^2 \mathbb{E}[\|\Gamma_{t-1}-\Gamma_t\|] \\
&\leq -\frac{\eta}{2}\|\Gamma_{t-1}\|_{\min}\,\mathbb{E}\|\nabla f(x_t)\|^2
   - \frac{\eta}{2}\mathbb{E}\!\left[\left\|\frac{1}{M}\sum_{m=1}^{M}\nabla f_m(x_t^m)\right\|_{\Gamma_{t-1}}^2\right] \\
&\quad + \frac{\eta}{2}\|\Gamma_{t-1}\|_{\max}\,\mathbb{E}\!\left[\left\|\frac{1}{M}\sum_{m=1}^{M}\nabla f_m(x_t) - \nabla f_m(x_t^m)\right\|^2\right]
   + \eta G^2 \mathbb{E}[\|\Gamma_{t-1}-\Gamma_t\|] \\
&\leq -\frac{\eta}{2C_0}\mathbb{E}\|\nabla f(x_t)\|^2
   - \frac{\eta}{2}\mathbb{E}\!\left[\left\|\frac{1}{M}\sum_{m=1}^{M}\nabla f_m(x_t^m)\right\|_{\Gamma_{t-1}}^2\right] \\
&\quad + \frac{\eta}{2\lambda M}\sum_{m=1}^{M}\mathbb{E}\!\left[\|\nabla f_m(x_t) - \nabla f_m(x_t^m)\|^2\right]
   + \eta G^2 \mathbb{E}[\|\Gamma_{t-1}-\Gamma_t\|] \\
&\leq -\frac{\eta}{2C_0}\mathbb{E}\|\nabla f(x_t)\|^2
   + \frac{\eta L^2}{2\lambda M}\sum_{m=1}^{M}\mathbb{E}\!\left[\|x_t - x_t^m\|^2\right]
   + \eta G^2 \mathbb{E}[\|\Gamma_{t-1}-\Gamma_t\|],
\end{align*}
where $\|\cdot\|$ indicates the spectral norm for matrices, and we used the following inequalities:
\[
\|\Gamma_{t-1}\|_{\min}
= \left\|\frac{1}{M}\sum_{m=1}^{M}\Gamma_{t-1}^m\right\|_{\min}
= \frac{1}{M}\sum_{m=1}^{M}\Gamma_{t-1}^m[i,i]
= \frac{1}{M}\sum_{m=1}^{M}\frac{1}{\sqrt{\hat{v}_{t-1}[i]+\lambda^2}}
\geq \frac{1}{\sqrt{G^2+\lambda^2}}
\stackrel{\text{def}}{=} \frac{1}{C_0}.
\]

%-----------------------------------------------------------------------------
% Step 3b (term II) — PDF page 38 top
%-----------------------------------------------------------------------------
\paragraph{Step 3b (one step progress). Bounding term II.}
\begin{align*}
II &= \eta\langle \nabla f(z_t), U_t \rangle
   \leq \eta\|\nabla f(z_t)\|\,\|U_t\|
   \leq \frac{\eta G}{M}\sum_{m=1}^{M}\|\Gamma_t^m(g_t - \widetilde{g}_t^m)\| \\
&\leq \frac{\eta G}{\lambda M}\sum_{m=1}^{M}\|g_t - \widetilde{g}_t^m\|.
\end{align*}

%-----------------------------------------------------------------------------
% Step 3c (term III) — PDF page 38
%-----------------------------------------------------------------------------
\paragraph{Step 3c (one step progress). Bounding term III.}
\begin{align*}
III &= \eta\langle \nabla f(z_t), V_t \rangle
   \leq \eta\|\nabla f(z_t)\|\,\|V_t\|
   \leq \frac{\eta\beta_1}{1-\beta_1}\frac{G}{M}\sum_{m=1}^{M}\|(\Gamma_{t-1}^m - \Gamma_t^m) u_{t-1}^m\| \\
&\leq \frac{\eta\beta_1}{1-\beta_1}\frac{G^2}{M}\sum_{m=1}^{M}\|\Gamma_{t-1}^m - \Gamma_t^m\|.
\end{align*}

%-----------------------------------------------------------------------------
% Step 3d (term IV) — PDF page 38
%-----------------------------------------------------------------------------
\paragraph{Step 3d (one step progress). Bounding term IV.}
\begin{align*}
IV &= \frac{\eta^2 L}{2}\|\Gamma_t g_t - U_t - V_t\|^2 \\
&\leq \frac{3\eta^2 L}{2}\|\Gamma_t g_t\|^2
   + \frac{3\eta^2 L}{2}\|U_t\|^2
   + \frac{3\eta^2 L}{2}\|V_t\|^2 \\
&\leq \frac{3\eta^2 L G^2}{2\lambda^2}
   + \frac{3\eta^2 L}{2\lambda^2 M}\sum_{m=1}^{M}\|g_t - \widetilde{g}_t^m\|^2
   + \frac{3\eta^2 \beta_1 L G}{2(1-\beta_1)M}\sum_{m=1}^{M}\|\Gamma_{t-1}^m - \Gamma_t^m\|^2
\end{align*}

%-----------------------------------------------------------------------------
% Step 3e (term V) — PDF page 38
%-----------------------------------------------------------------------------
\paragraph{Step 3e (one step progress). Bounding term V.}
\begin{align*}
V &= \eta\langle \nabla f(x_t) - \nabla f(z_t),\, \Gamma_t g_t \rangle \\
&= \eta\,\mathbb{E}\!\left[\langle \nabla f(x_t)-\nabla f(z_t),\, \Gamma_{t-1} g_t \rangle\right]
   + \eta\,\mathbb{E}\!\left[\langle \nabla f(x_t)-\nabla f(z_t),\, (\Gamma_t-\Gamma_{t-1}) g_t \rangle\right] \\
&\leq \eta\,\mathbb{E}\!\left[\left\langle \nabla f(x_t)-\nabla f(z_t),\, \frac{1}{M}\sum_{m=1}^{M}\nabla f_m(x_t^m) \right\rangle_{\!\Gamma_{t-1}}\right]
   + \frac{\eta^2 L \beta_1}{1-\beta_1}\mathbb{E}\!\left[\left\|\mathbb{E}_m[\Gamma_{t-1}^m u_{t-1}^m]\right\|\,\|(\Gamma_t-\Gamma_{t-1}) g_t\|\right] \\
&\leq \eta\,\mathbb{E}\!\left[\langle \nabla f(x_t)-\nabla f(z_t),\, \nabla f(x_t) \rangle_{\Gamma_{t-1}}\right] \\
&\quad + \eta\,\mathbb{E}\!\left[\left\langle \nabla f(x_t)-\nabla f(z_t),\, \frac{1}{M}\sum_{m=1}^{M}\nabla f_m(x_t^m)-\nabla f_m(x_t) \right\rangle_{\!\Gamma_{t-1}}\right]
   + \frac{\eta^2 L \beta_1 G^2}{(1-\beta_1)\lambda}\mathbb{E}\!\left[\|\Gamma_t-\Gamma_{t-1}\|\right] \\
&\leq \frac{\eta}{\lambda}\mathbb{E}\!\left[\|\nabla f(x_t)-\nabla f(z_t)\|\,\|\nabla f(x_t)\|\right] \\
&\quad + \frac{\eta}{\lambda}\mathbb{E}\!\left[\|\nabla f(x_t)-\nabla f(z_t)\|\cdot\frac{1}{M}\sum_{m=1}^{M}\|\nabla f_m(x_t^m)-\nabla f_m(x_t)\|\right]
   + \frac{\eta^2 L \beta_1 G^2}{(1-\beta_1)\lambda}\mathbb{E}\!\left[\|\Gamma_t-\Gamma_{t-1}\|\right] \\
&\leq \frac{\eta}{\lambda}\mathbb{E}\!\left[\frac{1}{2\rho}\|\nabla f(x_t)-\nabla f(z_t)\|^2 + \frac{\rho}{2}\|\nabla f(x_t)\|^2\right] \\
&\quad + \frac{\eta}{\lambda}\mathbb{E}\!\left[\frac{1}{2}\|\nabla f(x_t)-\nabla f(z_t)\|^2 + \frac{1}{2}\frac{L^2}{M}\sum_{m=1}^{M}\|x_t^m - x_t\|^2\right]
   + \frac{\eta^2 L \beta_1 G^2}{(1-\beta_1)\lambda}\mathbb{E}\!\left[\|\Gamma_t-\Gamma_{t-1}\|\right],
\end{align*}
where we used the following uniform bound on $\|\nabla f(x_t)-\nabla f(z_t)\|$:
\begin{align*}
\|\nabla f(x_t)-\nabla f(z_t)\|
&\leq L\|x_t - z_t\|
\leq \frac{\beta_1 L}{1-\beta_1}\|x_t - x_{t-1}\|
= \frac{\eta\beta_1 L}{1-\beta_1}\left\|\mathbb{E}_m[\Gamma_{t-1}^m u_{t-1}^m]\right\| \\
&\leq \frac{\eta\beta_1 L}{1-\beta_1}\mathbb{E}_m[\|\Gamma_{t-1}^m\|\,\|u_{t-1}^m\|]
\leq \frac{\eta\beta_1 L}{1-\beta_1}\frac{G}{\lambda}.
\end{align*}

%-----------------------------------------------------------------------------
% Summary bounds (O(·) simplified forms) — PDF page 39 top
%-----------------------------------------------------------------------------
Therefore, ignoring the constants, we have the following bounds:
\begin{align*}
V  &\leq \mathcal{O}\!\left(\frac{\eta^2}{\rho}\right)
   + \frac{\eta\rho}{2\lambda}\cdot\mathbb{E}[\|\nabla f(x_t)\|^2]
   + \mathcal{O}(\eta)\cdot\frac{1}{M}\sum_{m=1}^{M}\mathbb{E}[\|x_t^m - x_t\|^2]
   + \mathcal{O}(\eta^2) \\
IV &\leq \mathcal{O}(\eta^2) \\
III &\leq \mathcal{O}(\eta)\cdot\frac{1}{M}\sum_{m=1}^{M}\mathbb{E}[\|\Gamma_{t-1}^m - \Gamma_t^m\|] \\
II &\leq \mathcal{O}(\eta)\cdot\frac{1}{M}\sum_{m=1}^{M}\mathbb{E}[\|g_t - \widetilde{g}_t^m\|] \\
I  &\leq -\frac{\eta}{2C_0}\mathbb{E}\|\nabla f(x_t)\|^2
   + \mathcal{O}(\eta)\cdot\frac{1}{M}\sum_{m=1}^{M}\mathbb{E}\!\left[\|x_t - x_t^m\|^2\right]
   + \mathcal{O}(\eta)\cdot\frac{1}{M}\sum_{m=1}^{M}\mathbb{E}[\|\Gamma_{t-1}^m - \Gamma_t^m\|]
\end{align*}

%-----------------------------------------------------------------------------
% Three O(1/√T) conditions + double geometric sum — PDF page 39
%-----------------------------------------------------------------------------
To get the $\mathcal{O}\!\left(\frac{1}{\sqrt{T}}\right)$ bound for the averaged gradients $\mathbb{E}[\|\nabla f(x_t)\|^2]$, note that we are left to choose small value for $\rho = \frac{\lambda}{2C_0}$ and show the following bounds:
\begin{align*}
\frac{1}{TM}\sum_{t=0}^{T-1}\sum_{m=1}^{M}\mathbb{E}[\|x_t^m - x_t\|^2] &= \mathcal{O}(\eta^2), &\qquad&\text{(extension of Lemma~\ref{lem:bound})} \\
\sum_{t=0}^{T-1}\mathbb{E}[\|\Gamma_{t-1}^m - \Gamma_t^m\|] &= \mathcal{O}(1), &\qquad&\text{(follows from AMSGrad normalization)} \\
\frac{1}{M}\sum_{t=0}^{T-1}\sum_{m=1}^{M}\mathbb{E}[\|g_t - \widetilde{g}_t^m\|] &= \mathcal{O}(1), &\qquad&\text{(see below).}
\end{align*}

For the last bound, we can use similar steps as in Lemma~\ref{lem:bound}, namely
\begin{align*}
\mathbb{E}[\|u_t - u_t^m\|]
&= p_u \cdot 0 + (1-p_u)\,\mathbb{E}[\|\beta_1 u_{t-1} + (1-\beta_1)g_t - (\beta_1 u_{t-1}^m + (1-\beta_1)g_t^m)\|] \\
&\leq (1-p_u)\beta_1\,\mathbb{E}[\|u_{t-1}-u_{t-1}^m\|] + (1-p_u)(1-\beta_1)\,\mathbb{E}[\|g_t - g_t^m\|] \\
&\leq (1-p_u)(1-\beta_1)\sum_{\tau=0}^{t}((1-p_u)\beta_1)^{t-\tau}\,\mathbb{E}[\|g_\tau - g_\tau^m\|].
\end{align*}
\begin{align*}
\mathbb{E}[\|g_t - \widetilde{g}_t^m\|]
&= \mathbb{E}\left\|\frac{u_t - \beta_1 u_{t-1}}{1-\beta_1} - \frac{u_t^m - \beta_1 u_{t-1}^m}{1-\beta_1}\right\| \\
&\leq \frac{\beta_1}{1-\beta_1}\mathbb{E}\|u_{t-1}-u_{t-1}^m\| + \frac{1}{1-\beta_1}\mathbb{E}[\|u_t - u_t^m\|] \\
&= \frac{1}{1-\beta_1}\sum_{\tau=t-1}^{t}\beta_1^{t-\tau}\,\mathbb{E}\|u_\tau - u_\tau^m\| \\
&= (1-p_u)\sum_{\tau=t-1}^{t}\sum_{\nu=0}^{\tau}\beta_1^{t-\tau}\,((1-p_u)\beta_1)^{\tau-\nu}\,\mathbb{E}[\|g_\nu - g_\nu^m\|] \\
&= \sum_{\tau=t}^{t+1}\sum_{\nu=0}^{\tau-1}\beta_1^{t-\tau}\,(\underbrace{(1-p_u)\beta_1}_{=\,q_2})^{\tau-\nu}\,\mathbb{E}[\|g_\nu - g_\nu^m\|],
\end{align*}
which has the same double geometric sum structure as~\eqref{eq:double_geom}.

%=============================================================================
% D.2 Key Lemmas — PDF pages 39-43
%=============================================================================
\subsection{Key Lemmas}
\label{sec:key_lemmas}

\begin{lemma}
\label{lem:virt_drift}
For all $T \geq 1$, we have
\begin{equation}
\label{eq:lem2}
\sum_{t=0}^{T-1}\|z_t - x_t\|^2
\leq \frac{\eta^2\beta^2}{(1-\beta)^2 M}\,T\sigma^2
   + \frac{\eta^2\beta^2}{(1-\beta)^2}\sum_{t=0}^{T-1}\mathbb{E}\left\|\frac{1}{M}\sum_{m=1}^{M}\nabla f_m(x_t^m)\right\|^2.
\end{equation}
\end{lemma}

\begin{proof}
Since $u_{-1}=0$, unrolling the update rule of momentum, for any $t \geq 0$ we get
\[
u_t = \beta u_{t-1} + (1-\beta)g_t = (1-\beta)\sum_{\tau=0}^{t}\beta^{t-\tau} g^\tau.
\]
Using this and the definition of the average iterates, we have
\[
z_t - x_t = \frac{\beta}{1-\beta}(x_t - x_{t-1}) = -\frac{\beta\eta}{1-\beta}u_t = -\beta\eta\sum_{\tau=0}^{t}\beta^{t-\tau} g_\tau.
\]
Using convexity of squared norm function and letting $s_t \stackrel{\text{def}}{=} \sum_{\tau=0}^{t}\beta^{t-\tau} = \frac{1-\beta^{t+1}}{1-\beta}$, for all $t \geq 0$, we have
\[
\|z_t - x_t\|^2
= \eta^2\beta^2 s_t^2 \left\|\sum_{\tau=0}^{t}\frac{\beta^{t-\tau}}{s_t} g_\tau\right\|^2
\leq \eta^2\beta^2 s_t^2 \sum_{\tau=0}^{t}\frac{\beta^{t-\tau}}{s_t}\|g_\tau\|^2
\leq \frac{\eta^2\beta^2}{1-\beta}\sum_{\tau=0}^{t}\beta^{t-\tau}\|g_\tau\|^2.
\]
Summing over the iterates yields
\begin{align*}
\sum_{t=0}^{T-1}\mathbb{E}\|z_t - x_t\|^2
&\leq \frac{\eta^2\beta^2}{1-\beta}\sum_{t=0}^{T-1}\sum_{\tau=0}^{t}\beta^{t-\tau}\,\mathbb{E}\|g_\tau\|^2 \\
&= \frac{\eta^2\beta^2}{1-\beta}\sum_{\tau=0}^{T-1}\sum_{t=\tau}^{T-1}\beta^{t-\tau}\,\mathbb{E}\|g_\tau\|^2 \\
&= \frac{\eta^2\beta^2}{1-\beta}\sum_{\tau=0}^{T-1}\frac{1-\beta^{T-\tau}}{1-\beta}\,\mathbb{E}\|g_\tau\|^2 \\
&\leq \frac{\eta^2\beta^2}{(1-\beta)^2}\sum_{\tau=0}^{T-1}\mathbb{E}\|g_\tau\|^2 \\
&= \frac{\eta^2\beta^2}{(1-\beta)^2}\sum_{\tau=0}^{T-1}\mathbb{E}\left\|\frac{1}{M}\sum_{m=1}^{M}g_\tau^m - \nabla f_m(x_\tau^m)\right\|^2
   + \frac{\eta^2\beta^2}{(1-\beta)^2}\sum_{\tau=0}^{T-1}\mathbb{E}\left\|\frac{1}{M}\sum_{m=1}^{M}\nabla f_m(x_\tau^m)\right\|^2 \\
&= \frac{\eta^2\beta^2}{(1-\beta)^2 M^2}\sum_{\tau=0}^{T-1}\sum_{m=1}^{M}\mathbb{E}\left\|g_\tau^m - \nabla f_m(x_\tau^m)\right\|^2
   + \frac{\eta^2\beta^2}{(1-\beta)^2}\sum_{\tau=0}^{T-1}\mathbb{E}\left\|\frac{1}{M}\sum_{m=1}^{M}\nabla f_m(x_\tau^m)\right\|^2 \\
&= \frac{\eta^2\beta^2}{(1-\beta)^2 M}\,T\sigma^2
   + \frac{\eta^2\beta^2}{(1-\beta)^2}\sum_{\tau=0}^{T-1}\mathbb{E}\left\|\frac{1}{M}\sum_{m=1}^{M}\nabla f_m(x_\tau^m)\right\|^2.
\end{align*}
\end{proof}

\begin{lemma}
\label{lem:bound}
If $24\eta^2 L^2 \psi \leq 1$, then
\[
\frac{1}{MT}\sum_{t=0}^{T-1}\sum_{m=1}^{M}\mathbb{E}\|x_t - x_t^m\|^2
\leq 12\eta^2(B^2-1)\psi \cdot \frac{1}{T}\sum_{t=0}^{T-1}\mathbb{E}\|\nabla f(x_t)\|^2
   + 4\eta^2\psi(\sigma^2 + 3G^2),
\]
where
\begin{equation}
\label{eq:double_geom}
\psi = \frac{4(1-p_x)}{p_x^2}\cdot\frac{(1-\beta)(1-p_u)}{1-(1-p_u)\beta}
\end{equation}
\end{lemma}

\begin{proof}
Let us expand the term $\mathbb{E}\|x_{t+1}-x_{t+1}^m\|^2$ using $x_{t+1}^m$'s probabilistic update rule:
\begin{align*}
\mathbb{E}\|x_{t+1}-x_{t+1}^m\|^2
&= p_x \cdot 0 + (1-p_x)\cdot\mathbb{E}\|x_t - \eta u_t - (x_t^m - \eta u_t^m)\|^2 \\
&= (1-p_x)\cdot\mathbb{E}\|x_t - x_t^m - \eta(u_t - u_t^m)\|^2 \\
&\leq (1-p_x)(1+s)\,\mathbb{E}\|x_t - x_t^m\|^2 + \eta^2(1-p_x)(1+\nicefrac{1}{s})\,\mathbb{E}\|u_t - u_t^m\|^2 \\
&\leq \eta^2(1-p_x)(1+\nicefrac{1}{s})\sum_{\tau=1}^{t}((1-p_x)(1+s))^{t-\tau}\,\mathbb{E}\|u_\tau - u_\tau^m\|^2.
\end{align*}
where $s > 0$ will be chosen later. Next we expand the term $\mathbb{E}\|u_t - u_t^m\|^2$ using $u_t^m$'s probabilistic update rule:
\begin{align*}
\mathbb{E}\|u_t - u_t^m\|^2
&= p_u \cdot 0 + (1-p_u)\cdot\mathbb{E}\left\|\frac{1}{M}\sum_{m=1}^{M}(\beta u_{t-1}^m + (1-\beta)g_{t-1}^m) - (\beta u_{t-1}^m + (1-\beta)g_{t-1}^m)\right\|^2 \\
&= (1-p_u)\,\mathbb{E}\|\beta(u_{t-1}-u_{t-1}^m) + (1-\beta)(g_{t-1}-g_{t-1}^m)\|^2 \\
&\leq (1-p_u)\beta\,\mathbb{E}\|u_{t-1}-u_{t-1}^m\|^2 + (1-p_u)(1-\beta)\,\mathbb{E}\|g_{t-1}-g_{t-1}^m\|^2 \\
&\leq (1-p_u)(1-\beta)\sum_{\tau=0}^{t-1}((1-p_u)\beta)^{t-1-\tau}\,\mathbb{E}\|g_\tau - g_\tau^m\|^2 \\
&\leq \frac{1-\beta}{\beta}\sum_{\tau=0}^{t-1}((1-p_u)\beta)^{t-\tau}\,\mathbb{E}\|g_\tau - g_\tau^m\|^2
\end{align*}

Denote $q_1 = (1-p_x)(1+s)$ and $q_2 = (1-p_u)\beta$. Combining the previous two bounds, we get
\begin{align*}
\frac{1}{M}\sum_{m=1}^{M}\mathbb{E}\|x_t - x_t^m\|^2
&\leq \eta^2(1-p_x)(1+\nicefrac{1}{s})\sum_{\tau=1}^{t}((1-p_x)(1+s))^{t-\tau}\frac{1}{M}\sum_{m=1}^{M}\mathbb{E}\|u_\tau - u_\tau^m\|^2 \\
&\leq \eta^2(1-p_x)(1+\nicefrac{1}{s})\sum_{\tau=1}^{t}((1-p_x)(1+s))^{t-\tau}\frac{1}{M}\sum_{m=1}^{M}\left[\frac{1-\beta}{\beta}\sum_{\nu=0}^{\tau-1}((1-p_u)\beta)^{\tau-\nu}\,\mathbb{E}\|g_\nu - g_\nu^m\|^2\right] \\
&= \eta^2(1-p_x)(1+\nicefrac{1}{s})\frac{1-\beta}{\beta}\sum_{\tau=1}^{t}\sum_{\nu=0}^{\tau-1} q_1^{t-\tau}\,q_2^{\tau-\nu}\left[\frac{1}{M}\sum_{m=1}^{M}\mathbb{E}\|g_\nu - g_\nu^m\|^2\right] \\
&= \eta^2(1-p_x)(1+\nicefrac{1}{s})\frac{1-\beta}{\beta}\sum_{\nu=0}^{t-1} q_2\,\frac{q_1^{t-\nu}-q_2^{t-\nu}}{q_1-q_2}\left[\frac{1}{M}\sum_{m=1}^{M}\mathbb{E}\|g_\nu - g_\nu^m\|^2\right] \\
&= \eta^2 \underbrace{(1-p_x)(1+\nicefrac{1}{s})(1-\beta)(1-p_u)}_{\stackrel{\text{def}}{=}\,\phi}\sum_{\nu=0}^{t-1}\frac{q_1^{t-\nu}-q_2^{t-\nu}}{q_1-q_2}\left[\frac{1}{M}\sum_{m=1}^{M}\mathbb{E}\|g_\nu - g_\nu^m\|^2\right].
\end{align*}
Next, we bound the gradient term above.
\begin{align*}
\frac{1}{M}\sum_{m=1}^{M}\mathbb{E}\|g_t^m - g_t\|^2
&= \frac{1}{M}\sum_{m=1}^{M}\mathbb{E}\left\|g_t^m - \frac{1}{M}\sum_{i=1}^{M}g_t^i\right\|^2 \\
&\leq \frac{2}{M}\sum_{m=1}^{M}\mathbb{E}\left\|g_t^m - \nabla f_m(x_t^m) - \frac{1}{M}\sum_{m=1}^{M}(g_t^m - \nabla f_m(x_t^m))\right\|^2 \\
&\quad + \frac{2}{M}\sum_{m=1}^{M}\mathbb{E}\left\|\nabla f_m(x_t^m) - \frac{1}{M}\sum_{m=1}^{M}\nabla f_m(x_t^m)\right\|^2 \\
\text{(Lemma~\ref{lem:het})} \quad
&\leq \frac{2}{M}\sum_{m=1}^{M}\mathbb{E}\|g_t^m - \nabla f_m(x_t^m)\|^2
   - 2\,\mathbb{E}\left\|\frac{1}{M}\sum_{m=1}^{M}(g_t^m - \nabla f_m(x_t^m))\right\|^2 \\
&\quad + \frac{12L^2}{M}\sum_{m=1}^{M}\mathbb{E}\|x_t - x_t^m\|^2
   + 6(B^2-1)\,\mathbb{E}\|\nabla f(x_t)\|^2 + 6G^2 \\
&\leq 2\sigma^2 + \frac{12L^2}{M}\sum_{m=1}^{M}\mathbb{E}\|x_t - x_t^m\|^2
   + 6(B^2-1)\,\mathbb{E}\|\nabla f(x_t)\|^2 + 6G^2.
\end{align*}

Again, plugging this bound to the previous one, we get
\begin{align*}
\frac{1}{MT}\sum_{t=0}^{T-1}\sum_{m=1}^{M}\mathbb{E}\|x_t - x_t^m\|^2
&\leq \frac{1}{MT}\sum_{t=1}^{T}\sum_{m=1}^{M}\mathbb{E}\|x_t - x_t^m\|^2 \\
&\leq \frac{\eta^2\phi}{T}\sum_{t=1}^{T}\sum_{\tau=0}^{t-1}\frac{q_1^{t-\tau}-q_2^{t-\tau}}{q_1-q_2}\left[\frac{1}{M}\sum_{m=1}^{M}\mathbb{E}\|g_\tau - g_\tau^m\|^2\right] \\
&= \frac{\eta^2\phi}{T}\sum_{\tau=0}^{T-1}\frac{1}{q_1-q_2}\left(\frac{q_1(1-q_1^{T-\tau})}{1-q_1} - \frac{q_2(1-q_2^{T-\tau})}{1-q_2}\right)\left[\frac{1}{M}\sum_{m=1}^{M}\mathbb{E}\|g_\tau - g_\tau^m\|^2\right] \\
&\leq \frac{\eta^2\phi}{T}\sum_{\tau=0}^{T-1}\frac{1}{q_1-q_2}\left(\frac{q_1}{1-q_1} - \frac{q_2}{1-q_2}\right)\left[\frac{1}{M}\sum_{m=1}^{M}\mathbb{E}\|g_\tau - g_\tau^m\|^2\right] \\
&= \frac{\eta^2\phi}{(1-q_1)(1-q_2)}\frac{1}{T}\sum_{\tau=0}^{T-1}\left[\frac{1}{M}\sum_{m=1}^{M}\mathbb{E}\|g_\tau - g_\tau^m\|^2\right].
\end{align*}

Now, let us optimize the factor
\[
\frac{\phi}{(1-q_1)(1-q_2)}
= \frac{(1-p_x)(1+\nicefrac{1}{s})(1-\beta)(1-p_u)}{(1-(1-p_x)(1+s))(1-(1-p_u)\beta)}
= \frac{(1-p_x)(1+\nicefrac{1}{s})}{1-(1-p_x)(1+s)}\cdot\frac{(1-\beta)(1-p_u)}{1-(1-p_u)\beta}
\]
by choosing optimal value for $s$ introduced earlier. By the first order optimality condition, we find that the optimal value is $s^* = \frac{1}{\sqrt{1-p_x}} - 1$. Hence, the minimal value of the factor is
\begin{align*}
\frac{\phi}{(1-q_1)(1-q_2)}
&= \frac{1-p_x}{(1-\sqrt{1-p_x})^2}\cdot\frac{(1-\beta)(1-p_u)}{1-(1-p_u)\beta} \\
&= \frac{(1-p_x)(1+\sqrt{1-p_x})^2}{p_x^2}\cdot\frac{(1-\beta)(1-p_u)}{1-(1-p_u)\beta} \\
&\leq \frac{4(1-p_x)}{p_x^2}\cdot\frac{(1-\beta)(1-p_u)}{1-(1-p_u)\beta}
\stackrel{\text{def}}{=} \psi.
\end{align*}

Continuing the chain of bounds
\begin{align*}
\frac{1}{MT}\sum_{t=0}^{T-1}\sum_{m=1}^{M}\mathbb{E}\|x_t - x_t^m\|^2
&\leq \eta^2\psi\cdot\frac{1}{T}\sum_{t=0}^{T-1}\left[\frac{1}{M}\sum_{m=1}^{M}\mathbb{E}\|g_t - g_t^m\|^2\right] \\
&\leq \eta^2\psi\cdot\frac{1}{T}\sum_{t=0}^{T-1}\left[\frac{12L^2}{M}\sum_{m=1}^{M}\mathbb{E}\|x_t - x_t^m\|^2 + 6(B^2-1)\,\mathbb{E}\|\nabla f(x_t)\|^2 + 2\sigma^2 + 6G^2\right] \\
&\leq 12\eta^2 L^2\psi\cdot\frac{1}{TM}\sum_{t=0}^{T-1}\sum_{m=1}^{M}\mathbb{E}\|x_t - x_t^m\|^2 \\
&\quad + 6\eta^2(B^2-1)\psi\cdot\frac{1}{T}\sum_{t=0}^{T-1}\mathbb{E}\|\nabla f(x_t)\|^2 + 2\eta^2\psi(\sigma^2 + 3G^2).
\end{align*}

Assuming $12\eta^2 L^2 \psi \leq \nicefrac{1}{2}$ and reordering the first term in the bound, we arrive
\[
\frac{1}{MT}\sum_{t=0}^{T-1}\sum_{m=1}^{M}\mathbb{E}\|x_t - x_t^m\|^2
\leq 12\eta^2(B^2-1)\psi\cdot\frac{1}{T}\sum_{t=0}^{T-1}\mathbb{E}\|\nabla f(x_t)\|^2 + 4\eta^2\psi(\sigma^2 + 3G^2).
\]
\end{proof}

\begin{lemma}
\label{lem:het}
Under smoothness and bounded heterogeneity assumptions~1 and~3, we have
\[
\frac{1}{M}\sum_{m=1}^{M}\left\|\nabla f_m(x_t^m) - \frac{1}{M}\sum_{i=1}^{M}\nabla f_i(x_t^i)\right\|^2
\leq \frac{6L^2}{M}\sum_{m=1}^{M}\|x_t - x_t^m\|^2 + 3(B^2-1)\|\nabla f(x_t)\|^2 + 3G^2.
\]
\end{lemma}

\begin{proof}
The bound follows from simple algebraic manipulations and Jensen's inequality.
\begin{align*}
&\frac{1}{M}\sum_{m=1}^{M}\left\|\nabla f_m(x_t^m) - \frac{1}{M}\sum_{i=1}^{M}\nabla f_i(x_t^i)\right\|^2 \\
&= \frac{1}{M}\sum_{m=1}^{M}\left\|\nabla f_m(x_t^m) - \nabla f_m(x_t) + \nabla f_m(x_t) - \nabla f(x_t) + \nabla f(x_t) - \frac{1}{M}\sum_{i=1}^{M}\nabla f_i(x_t^i)\right\|^2 \\
&\leq \frac{3}{M}\sum_{m=1}^{M}\|\nabla f_m(x_t^m) - \nabla f_m(x_t)\|^2
   + \frac{3}{M}\sum_{m=1}^{M}\|\nabla f_m(x_t) - \nabla f(x_t)\|^2 \\
&\quad + \frac{3}{M}\sum_{m=1}^{M}\left\|\nabla f(x_t) - \frac{1}{M}\sum_{i=1}^{M}\nabla f_i(x_t^i)\right\|^2 \\
&\leq \frac{3L^2}{M}\sum_{m=1}^{M}\|x_t^m - x_t\|^2
   + \frac{3}{M}\sum_{m=1}^{M}\|\nabla f_m(x_t) - \nabla f(x_t)\|^2
   + \frac{3L^2}{M}\sum_{i=1}^{M}\|x_t - x_t^i\|^2 \\
&= \frac{6L^2}{M}\sum_{m=1}^{M}\|x_t^m - x_t\|^2
   + \frac{3}{M}\sum_{m=1}^{M}\|\nabla f_m(x_t) - \nabla f(x_t)\|^2 \\
&= \frac{6L^2}{M}\sum_{m=1}^{M}\|x_t^m - x_t\|^2 + 3G^2 + 3(B^2-1)\|\nabla f(x_t)\|^2.
\end{align*}
\end{proof}

%=============================================================================
% Appendix E: High-probability bounds for DES-LOC-Adam
% NOTE: Full proof deferred to Claude instance 14 (M326-M350)
%=============================================================================
\section{Convergence Analysis of DES-LOC-Adam (high-probability bounds)}
\label{sec:high_prob}
% Placeholder — full proof spans pages 44-46 of the PDF

%=============================================================================
\section{Derivation of Eqs.\ (1) and (2): Maximum Momentum Change With Clipping}
\label{sec:drift_derivation}
%=============================================================================

\textbf{Step 2: Bound on $\|u_{t+K} - u_t\|_\infty$.} Now we bound the change in the momentum over $K$ steps
explicitly. Unrolling the recursion, we have:
\begin{equation}
u_{t+K} = \beta_1^K u_t + (1 - \beta_1) \sum_{k=0}^{K-1} \beta_1^k g_{t+K-k}.
\label{eq:unroll_u}
\end{equation}

Subtracting $u_t$ from both sides, we obtain:
\begin{equation}
u_{t+K} - u_t = (\beta_1^K - 1)u_t + (1 - \beta_1) \sum_{k=0}^{K-1} \beta_1^k g_{t+K-k}.
\label{eq:diff_u}
\end{equation}

Applying the triangle inequality gives:
\begin{equation}
\|u_{t+K} - u_t\|_\infty \leq |1 - \beta_1^K| \|u_t\|_\infty + (1 - \beta_1) \sum_{k=0}^{K-1} \beta_1^k \|g_{t+K-k}\|_\infty.
\label{eq:triangle_u}
\end{equation}

Using $\|g_{t,i}\|_\infty \leq \rho$ (from coordinate-wise clipping) and $\|u_t\|_\infty \leq \rho$ (by induction, since $u$ is an exponential moving average of clipped gradients), the geometric series evaluates to:
\begin{equation}
\|u_{t+K} - u_t\|_\infty \leq (1 - \beta_1^K)\rho + (1 - \beta_1^K)\rho = 2\rho(1 - \beta_1^K).
\end{equation}

The identical argument applies to the second moment $v_t$ with $\rho$ replaced by $\rho^2$ (since $(g_t \odot g_t)_i \leq \rho^2$), yielding:
\begin{equation}
\|v_{t+K} - v_t\|_\infty \leq 2\rho^2(1 - \beta_2^K).
\end{equation}

%=============================================================================
\section{Wall-Clock Time Modeling}
\label{sec:wallclock_model}
%=============================================================================

\subsection{Estimating Total Wall-Clock Time}

We model the total wall-clock time $\mathcal{T}$ for training as:
\begin{equation}
\mathcal{T} = T \cdot t_{\text{step}} + R_x \cdot t_{\text{comm}}(x) + \sum_{j=1}^{N} R_j \cdot t_{\text{comm}}(s^j),
\end{equation}
where $T$ is the total number of steps, $t_{\text{step}}$ is the per-step computation time, $R_x = T/K_x$ is the number of parameter sync rounds, $R_j = T/K_j$ is the number of sync rounds for state $j$, and $t_{\text{comm}}(\cdot)$ is the communication time for a given tensor.

For Ring-AllReduce with model size $|x|$ parameters, bandwidth $B_w$ (bytes/sec), and $M$ workers:
\begin{equation}
t_{\text{comm}}(x) = \frac{2(M-1)}{M} \cdot \frac{|x| \cdot \text{bytes\_per\_param}}{B_w}.
\end{equation}

DES-LOC's communication cost relative to DDP is:
\begin{equation}
\frac{\mathcal{T}_{\text{DES-LOC}}}{\mathcal{T}_{\text{DDP}}} = \frac{T \cdot t_{\text{step}} + \frac{T}{K_x} t_{\text{comm}}(x) + \frac{T}{K_u} t_{\text{comm}}(u) + \frac{T}{K_v} t_{\text{comm}}(v)}{T \cdot (t_{\text{step}} + t_{\text{comm}}(x) + t_{\text{comm}}(u) + t_{\text{comm}}(v))}.
\end{equation}

With $K_x = 256$, $K_u = 3K_x = 768$, $K_v = 6K_x = 1536$, DES-LOC communicates $\approx 170\times$ less than DDP.

\begin{table}[h]
\centering
\caption{Throughput (tokens/sec/GPU) at different scales. DES-LOC matches FAVG+OPT throughput.}
\label{tab:throughput}
\begin{tabular}{l|cc|cc}
\toprule
& \multicolumn{2}{c|}{1B Model} & \multicolumn{2}{c}{7B Model} \\
$K_x$ & 16 & 256 & 16 & 256 \\
\midrule
DDP          & 60k & 60k & 8.5k & 8.5k \\
FAVG+OPT     & 106k & 135k & 11.5k & 13.7k \\
Local Adam   & 89k & 132k & 10.4k & 13.6k \\
DES-LOC      & 105k & 134k & 11.4k & 13.7k \\
\bottomrule
\end{tabular}
\end{table}

%=============================================================================
\section{Checkpointing vs.\ Periodic State Synchronization}
\label{sec:checkpoint}
%=============================================================================

A natural question is whether periodic checkpointing can replace DES-LOC's explicit state synchronization for handling worker failures. We argue it cannot, for two reasons:

\begin{enumerate}
\item \textbf{State staleness.} When a new worker joins from a checkpoint, its optimizer states ($u, v$) reflect the global state at checkpoint time. If the checkpoint is old, these states are stale and can cause training instability. DES-LOC's periodic synchronization ensures all workers' states stay within bounded drift (Eqs.~\ref{eq:drift_u},~\ref{eq:drift_v}).

\item \textbf{No convergence guarantees.} Ad-hoc checkpoint averaging provides no theoretical guarantees about the quality of the resulting optimizer states. DES-LOC's provable convergence (Theorem~\ref{thm:sgdm}) holds precisely because synchronization is periodic and bounded.
\end{enumerate}

%=============================================================================
\section{Choosing Synchronization Frequencies}
\label{sec:choosing_freq}
%=============================================================================

For a state with a decay rate $\beta$ approaching 1, the half-life can be calculated as:
\begin{equation}
t_{1/2} \approx \frac{\ln 2}{1 - \beta}.
\end{equation}

Based on this, we propose the following two-step methodology:

\paragraph{Parameters First ($K_x$).} The synchronization of model parameters is paramount to training
quality. The period $K_x$ should be chosen to match the end-of-training quality of fully-synchronous DDP at a target step budget while still materially reducing communication. In
practice, starting points like $K_x = 16$ or $K_x = 32$ are effective.

\paragraph{Momentum by Half-Life ($K_u, K_v$).} For any optimizer momentum state with a decay rate
$\beta$, its synchronization period $K$ should be set near its calculated half-life, i.e., $K \approx t_{1/2}$.
For common optimizers like Adam or ADOPT with well-tuned decay rates $\beta_1$ and $\beta_2$, this
simplifies to:
\begin{equation}
K_u \approx \frac{\ln 2}{1 - \beta_1}, \quad K_v \approx \frac{\ln 2}{1 - \beta_2}.
\end{equation}

Following this heuristic can yield a minimum $5\times$ reduction in communication cost over DDP for
$K_x \geq 16$, significantly decreasing wall-clock time while achieving convergence speed and final
model quality comparable to DDP.

%=============================================================================
\section{Critical Batch Size and Regime Positioning}
\label{sec:critical_batch}
%=============================================================================

To formally contextualize the regimes where DES-LOC is most beneficial, we begin with the statistical
properties of gradient estimation. Let the true gradient over the full data distribution for a loss function
$L: \mathbb{R}^d \to \mathbb{R}$ be $G(\theta) = \nabla L(\theta)$. In practice, a mini-batch of size $B$ provides an estimate, $G_{\text{est}}(\theta)$.
The variance of this estimator scales inversely with the batch size:
\begin{equation}
\text{Var}[G_{\text{est}}(\theta)] = \frac{1}{B}\Sigma(\theta),
\end{equation}
where $\Sigma(\theta)$ is the per-sample gradient covariance. DES-LOC is most beneficial in the large-batch regime where communication costs dominate compute costs.

%=============================================================================
\section{Extended Related Work}
\label{sec:extended_related}
%=============================================================================

\paragraph{Gradient compression and sparsification.}
Methods like Deep Gradient Compression~\citep{lin2018deep}, TopK sparsification~\citep{alistarh2018convergence}, and CocktailSGD~\citep{wang2023cocktailsgd} reduce communication by compressing or sparsifying gradients. These are orthogonal to DES-LOC and can be combined for further savings.

\paragraph{Model parallelism.}
Megatron-LM~\citep{shoeybi2019megatron}, ZeRO~\citep{rajbhandari2020zero}, and FSDP~\citep{zhao2023pytorch} partition model parameters across devices. DES-LOC focuses on data parallelism and is complementary to model-parallel approaches.

\paragraph{Personalized and heterogeneous federated learning.}
Methods like Ditto~\citep{li2021ditto}, FedProx~\citep{li2020federated}, and Per-FedAvg~\citep{fallah2020personalized} address non-IID data in federated settings. DES-LOC inherits robustness to heterogeneity from its theoretical framework (Assumption~\ref{asm:heterogeneity}).

%=============================================================================
\section{LLM Usage Declaration}
\label{sec:llm_usage}
%=============================================================================

Large Language Models were used during the preparation of this work for editing and improving the
text quality of the manuscript. No LLM was used for generating scientific content, proofs, or
experimental results. All technical contributions are the original work of the authors.

%=============================================================================
\section{Limitations}
\label{sec:limitations}
%=============================================================================

\begin{enumerate}
\item \textbf{Communication overlap.} Our current implementation uses a ``stop-the-world'' synchronization approach. Overlapping communication with computation~\citep{douillard2025streaming} could further improve wall-clock time.

\item \textbf{Heterogeneous hardware.} We assume homogeneous worker hardware. Extending DES-LOC to heterogeneous clusters with different compute speeds is left for future work.

\item \textbf{Adaptive synchronization.} Our synchronization periods are fixed throughout training. Adapting them based on training dynamics (e.g., increasing $K_x$ as training progresses) could yield further improvements.

\item \textbf{Scale of experiments.} While we validate up to 1.7B parameters and project to 13B, experiments at truly large scale (100B+) remain future work due to resource constraints.
\end{enumerate}

\end{document}
